% !TeX program = xelatex
% 数译（mathtranslations.org）数学中文译本排版模板（讲义 ctexart 版）
% Keith Conrad: Tensor Products (Part I)
%
\documentclass[UTF8,12pt,fontset=none]{ctexart}

% ---------- 页面、字体与数学 ----------
\usepackage[paper=a4paper,left=2.5cm,right=2.5cm,top=2.8cm,bottom=2.8cm]{geometry}
\usepackage{fontspec}
\usepackage{xeCJK}
\usepackage{amsmath,amsthm,amssymb,mathrsfs,mathtools}
\usepackage{graphicx,tikz,tikz-cd,titlesec}
\usepackage{array,booktabs,longtable}
\usepackage{etoolbox}
\usepackage[dvipsnames]{xcolor}
\usepackage[
    colorlinks=true,
    bookmarksnumbered=true,
    linkcolor=MidnightBlue,
    citecolor=BrickRed,
    urlcolor=MidnightBlue,
    linktoc=all
]{hyperref}

% 允许版心稍微拉伸以避免孤立过宽行
\setlength{\emergencystretch}{3em}

% 英文：Computer Modern Unicode Serif
\setmainfont{cmunrm.otf}[
  BoldFont=cmunbx.otf,
  ItalicFont=cmunti.otf,
  BoldItalicFont=cmunbi.otf
]

% 中文：正文宋体；专用术语楷体；定理类环境仿宋
\setCJKmainfont{FandolSong-Regular.otf}[BoldFont=FandolSong-Bold.otf]
\newCJKfontfamily\songfont{FandolSong-Regular.otf}[BoldFont=FandolSong-Bold.otf]
\newCJKfontfamily\termfont{FandolKai-Regular.otf}
\newCJKfontfamily\theoremfont[AutoFakeBold=2]{FandolFang-Regular.otf}
\setCJKsansfont{FandolHei-Regular.otf}[BoldFont=FandolHei-Bold.otf]
\setCJKmonofont{FandolFang-Regular.otf}

\pagestyle{plain}
\setlength{\parindent}{2em}
\setlength{\parskip}{0pt}
\linespread{1.16}

% ---------- 译本元信息 ----------
\newcommand{\BookTitleCN}{张量积}
\newcommand{\BookTitleEN}{Tensor Products}
\newcommand{\OriginalAuthor}{Keith Conrad}
\newcommand{\OriginalEdition}{Expository Notes}
\newcommand{\OriginalPublisher}{University of Connecticut}
\newcommand{\OriginalYear}{2024}
\newcommand{\Translator}{数译}
\newcommand{\ModelUsed}{GLM-5 / Codex}
\newcommand{\TranslationDate}{\today}

% ================================================================
%                         术语系统
% ================================================================
\newcommand{\termcn}[1]{{\termfont #1}}

\makeatletter
\newcommand{\MTTermList}{}
\newcommand{\MTRegisterTerm}[3]{%
  \g@addto@macro\MTTermList{\MTTermRow{#1}{#2}{#3}}%
}
\newcommand{\MTTermRow}[3]{%
  {\termfont #1} & #2 & \hyperref[#3]{\pageref*{#3}} \\
}

\newcommand{\newterm}[3]{%
  \@ifundefined{MT@seen@#1}{%
    \expandafter\gdef\csname MT@seen@#1\endcsname{}%
    \phantomsection
    \label{term:#1}%
    \protected@write\@auxout{}{%
      \string\MTRegisterTerm{\unexpanded{#2}}{\unexpanded{#3}}{term:#1}%
    }%
  }{}%
  {\termfont #2}（#3）%
}
\makeatother

\newcommand{\printterminology}{%
  \clearpage
  \phantomsection
  \section*{术语索引}
  \addcontentsline{toc}{section}{术语索引}

  \vspace{0.5cm}
  \begingroup
  \renewcommand{\arraystretch}{1.28}
  \begin{longtable}{>{\raggedright\arraybackslash}p{0.28\textwidth}>{\raggedright\arraybackslash}p{0.49\textwidth}>{\centering\arraybackslash}p{0.12\textwidth}}
    \toprule
    \textbf{中文术语} & \textbf{英文术语} & \textbf{页码} \\
    \midrule
    \endfirsthead
    \toprule
    \textbf{中文术语} & \textbf{英文术语} & \textbf{页码} \\
    \midrule
    \endhead
    \MTTermList
    \bottomrule
  \end{longtable}
  \endgroup
}

% ================================================================
%                       定理类环境
% ================================================================
\newtheoremstyle{mtfang}%
  {6pt}{6pt}%
  {\theoremfont}%
  {0pt}%
  {\songfont\bfseries}%
  {：}%
  {0.5em}%
  {}%

\newtheoremstyle{mtsong}%
  {6pt}{6pt}%
  {\songfont}%
  {0pt}%
  {\songfont\bfseries}%
  {：}%
  {0.5em}%
  {}%

% 讲义中公式随节编号
\numberwithin{equation}{section}

% 所有定理类环境共用节级计数器，与原书完全一致
\theoremstyle{mtfang}
\newtheorem{theorem}{定理}[section]
\newtheorem{lemma}[theorem]{引理}
\newtheorem{proposition}[theorem]{命题}
\newtheorem{corollary}[theorem]{推论}
\newtheorem{definition}[theorem]{定义}
\newtheorem{remark}[theorem]{注}

\theoremstyle{mtsong}
\newtheorem{example}[theorem]{例}
\newtheorem{problem}[theorem]{问题}
\newtheorem{question}[theorem]{问题}

\newcommand{\envenddiamond}{%
  \ifhmode
    \unskip\nobreak\hfill\ensuremath{\diamond}%
  \else
    \noindent\hfill\ensuremath{\diamond}%
  \fi
  \par
}
\AtEndEnvironment{example}{\envenddiamond}
\AtEndEnvironment{remark}{\envenddiamond}

% proof 样式
\makeatletter
\renewcommand{\proofname}{证明}
\renewenvironment{proof}[1][\proofname]{\par
  \pushQED{\qed}%
  \normalfont\songfont
  \topsep6\p@\@plus6\p@\relax
  \trivlist
  \item[\hskip\labelsep\normalfont\songfont\bfseries #1：]\ignorespaces
}{%
  \popQED\endtrivlist\@endpefalse
}
\makeatother
\renewcommand{\qedsymbol}{\ensuremath{\square}}

% ================================================================
%                      讲义节标题格式
% ================================================================
\titleformat{\section}[block]
  {\normalfont\Large\bfseries}
  {\S\,\thesection\quad}
  {0pt}
  {}
\titlespacing*{\section}{0pt}{2ex plus 0.5ex minus .2ex}{1.2ex plus .2ex}

\ctexset{
  subsection={
    format=\large\bfseries,
    name={,.\quad},
    number=\arabic{section}.\arabic{subsection}
  }
}

% 常用数学记号
\newcommand{\R}{\mathbf{R}}
\newcommand{\C}{\mathbf{C}}
\newcommand{\Z}{\mathbf{Z}}
\newcommand{\Q}{\mathbf{Q}}
\newcommand{\N}{\mathbf{N}}
\newcommand{\F}{\mathbf{F}}
\newcommand{\Hom}{\operatorname{Hom}}
\newcommand{\End}{\operatorname{End}}
\newcommand{\Aut}{\operatorname{Aut}}
\newcommand{\im}{\operatorname{im}}
\newcommand{\id}{\operatorname{id}}
\newcommand{\coker}{\operatorname{coker}}
\newcommand{\Tor}{\operatorname{Tor}}
\newcommand{\Ext}{\operatorname{Ext}}
\newcommand{\ann}{\operatorname{ann}}

% 参考文献环境
\renewcommand{\refname}{参考文献}

\begin{document}

% ================================================================
%                             封面
% ================================================================
\begin{titlepage}
    \thispagestyle{empty}
    \flushleft

    \rule{\textwidth}{1.2pt} \\[0.1cm]
    \rule{\textwidth}{0.8pt} \\[1.5cm]

    \hspace{2cm} {\fontsize{42pt}{52pt}\selectfont \bfseries \BookTitleCN} \\[0.8cm]
    \hspace{2cm} {\fontsize{24pt}{28pt}\selectfont \textit{\BookTitleEN}} \\[1.5cm]

    \rule{\textwidth}{0.6pt} \\[2cm]

    \hspace{2cm} {\Huge \textsc{\OriginalAuthor}} \\[5cm]

    \vfill

    \noindent
    \begin{minipage}[b]{0.46\textwidth}
        \raggedright
        \includegraphics[width=5.7cm]{logo.pdf}
    \end{minipage}
    \hfill
    \begin{minipage}[b]{0.50\textwidth}
        \raggedleft
        {\small \OriginalEdition · \OriginalPublisher · \OriginalYear} \\[0.25cm]
        {\normalsize \Translator\ 翻译及重排} \\[0.25cm]
        {\small Model: \ModelUsed} \\
        {\small Updated: \TranslationDate}
    \end{minipage}

    \vspace*{1cm}
\end{titlepage}

% ================================================================
%                       译者说明与目录
% ================================================================
\pagenumbering{roman}
\setcounter{page}{1}

\phantomsection
\section*{译者说明}
\addcontentsline{toc}{section}{译者说明}

本讲义是康涅狄格大学（University of Connecticut）Keith Conrad 教授所撰写的高等代数著名讲义《张量积》（\textit{Tensor Products}）前篇（Part I）的中译本.
Keith Conrad 的讲义以循序渐进、直观透彻与例题详实著称，是学习现代交换代数、模论与张量积普遍映射性质（Universal Mapping Property）的极佳教材.

本次翻译与重排严格遵循以下原则：
\begin{enumerate}
    \item \textbf{忠实原著}：以 Keith Conrad 原讲义 PDF 为唯一权威，完整保留原讲义的所有定理、引理、例题、注记、历史渊源与参考文献体系；
    \item \textbf{统一编号}：严格保持原讲义各定理类环境（定义、定理、例、注等）共用节级计数器的设计，确保中译本的所有引证编号与原讲义 1:1 精确对应；
    \item \textbf{现代排版与超链接}：所有普遍映射性质的交换图均使用 \texttt{tikz-cd} 矢量重绘；全文公式与定理引用均已实现可点击的双向超链接；
    \item \textbf{术语索引}：首次出现的专有名词使用 \texttt{\string\newterm} 标记，并在文末附有中英对照术语索引表.
\end{enumerate}

\vspace{0.5cm}
\hfill \Translator

\hfill \TranslationDate

\newpage
\begingroup
\hypersetup{linkcolor=black}
\tableofcontents
\endgroup
\clearpage

\pagenumbering{arabic}
\setcounter{page}{1}

% sections/sec1_intro.tex
\section{引言}
\label{sec:intro}

设 $R$ 为交换环，并设 $M$ 与 $N$ 为 $R$-模.（我们始终假定所考虑的环具有乘法幺元，且模均假定为幺模：对所有 $m \in M$ 均有 $1 \cdot m = m$.）直和 $M \oplus N$ 是模的一种加法运算.在此，我们引入一种乘法运算 $M \otimes_R N$，称为\newterm{tensor-product}{张量积}{tensor product}.我们首先描述模的张量积应当具有怎样的形式，严格的定义见第~\ref{sec:construction}~节.

张量积最早起源于向量空间，这也是物理学与工程学中唯一出现张量积的背景，因此我们首先叙述向量空间的张量积.设 $V$ 与 $W$ 为域 $K$ 上的向量空间，并分别选取 $V$ 的一组基 $\{e_i\}$ 与 $W$ 的一组基 $\{f_j\}$.张量积 $V \otimes_K W$ 定义为以形式符号 $e_i \otimes f_j$ 为基的 $K$-向量空间（我们在定义上宣告这些新符号是线性无关的）.因此，$V \otimes_K W$ 由形如 $\sum_{i,j} c_{ij} e_i \otimes f_j$ 的形式和组成，其中 $c_{ij} \in K$.$V \otimes_K W$ 中的元素称为\newterm{tensor}{张量}{tensor}.对于 $v \in V$ 与 $w \in W$，将 $v$ 与 $w$ 用 $V$ 和 $W$ 的原始基展开，然后将 $v \otimes w$ 展开，仿佛 $\otimes$ 是一种非交换乘法（且允许提出标量因子），以此方式得到的 $V \otimes_K W$ 中的元素定义为 $v \otimes w$.

例如，设 $V = W = \R^2 = \R e_1 + \R e_2$，其中 $\{e_1, e_2\}$ 是标准基.（对两个 $\R^2$ 我们使用相同的基.）则 $\R^2 \otimes_\R \R^2$ 是一个以 $e_1 \otimes e_1$、$e_1 \otimes e_2$、$e_2 \otimes e_1$ 与 $e_2 \otimes e_2$ 为基的 $4$ 维空间.若 $v = e_1 - e_2$ 且 $w = e_1 + 2e_2$，则
\begin{equation}
  v \otimes w = (e_1 - e_2) \otimes (e_1 + 2e_2) := e_1 \otimes e_1 + 2e_1 \otimes e_2 - e_2 \otimes e_1 - 2e_2 \otimes e_2. \label{eq:1.1}
\end{equation}

$v \otimes w$ 是否依赖于 $\R^2$ 基的选取？作为检验，另选一组基，例如 $e'_1 = e_1 + e_2$ 与 $e'_2 = 2e_1 - e_2$.则 $v$ 与 $w$ 可以表示为 $v = -\frac{1}{3}e'_1 + \frac{2}{3}e'_2$ 以及 $w = \frac{5}{3}e'_1 - \frac{1}{3}e'_2$.通过形式计算，
\begin{align*}
  v \otimes w &= \left(-\frac{1}{3}e'_1 + \frac{2}{3}e'_2\right) \otimes \left(\frac{5}{3}e'_1 - \frac{1}{3}e'_2\right) \\
  &= -\frac{5}{9}e'_1 \otimes e'_1 + \frac{1}{9}e'_1 \otimes e'_2 + \frac{10}{9}e'_2 \otimes e'_1 - \frac{2}{9}e'_2 \otimes e'_2,
\end{align*}
如果将这最后一个线性组合中 $e'_1$ 与 $e'_2$ 关于 $e_1$ 和 $e_2$ 的定义代入、完全展开并合并同类项，便会回到公式~\eqref{eq:1.1} 等号右侧的和式.这表明 $v \otimes w$ 在 $\R^2 \otimes_\R \R^2$ 中具有与基的选取无关的明确意义，尽管要一般性地证明这一点看似令人望而生畏.

在模的背景下，张量积可以类似向量空间的情形来描述，但 $\otimes$ 所应满足的性质必须在一般意义下规定，而不能仅仅依赖于基（因为模可能根本不存在基）：对于 $R$-模 $M$ 与 $N$，它们的张量积 $M \otimes_R N$（读作 ``$M$ tensor $N$'' 或 ``$M$ 与 $N$ 在 $R$ 上的张量积''）是一个由所有形如 $m \otimes n$（其中 $m \in M, n \in N$）的符号所生成的 $R$-模——这里不是作为一组基，而仅仅作为生成集\footnote{回忆一下，$R$-模的一个生成集是指其有限 $R$-线性组合能充满整个模的子集.生成集总是存在的，因为整个模本身就是一个生成集.}——并且这些符号满足分配律：
\begin{equation}
  (m + m') \otimes n = m \otimes n + m' \otimes n, \quad m \otimes (n + n') = m \otimes n + m \otimes n'. \label{eq:1.2}
\end{equation}
此外，乘以标量 $r \in R$ 可以放入 $\otimes$ 的任意一侧：对 $m \in M$ 与 $n \in N$，
\begin{equation}
  r(m \otimes n) = (rm) \otimes n = m \otimes (rn). \label{eq:1.3}
\end{equation}
因此，书写 $rm \otimes n$ 是毫无歧义的：它既表示 $r(m \otimes n)$，也表示 $(rm) \otimes n$.

应当将 $M \otimes_R N$ 中的公式~\eqref{eq:1.2} 与~\eqref{eq:1.3} 和直和 $M \oplus N$ 中的公式加以对比，在直和中：
\[
  (m + m', n) = (m, n) + (m', 0), \quad r(m, n) = (rm, rn).
\]
在 $M \oplus N$ 中，一个元素 $(m, n)$ 可分解为 $(m, 0) + (0, n)$，但 $M \otimes_R N$ 中的元素 $m \otimes n$ 却无法这样拆开.虽然 $M \oplus N$ 中的每个元素都是有序对 $(m, n)$，但在 $M \otimes_R N$ 中，通常存在远多于形如 $m \otimes n$ 的乘积元素.$M \otimes_R N$ 中的一般元素称为\termcn{张量}，它是有限个此类积的 $R$-线性组合\footnote{作为对比，考虑多项式环 $R[X, Y]$，其元素不仅包括乘积项 $f(X)g(Y)$，更包括这些乘积项的和式 $\sum_{i,j} a_{ij} X^i Y^j$.事实上，作为 $R$-模有 $R[X, Y] \cong R[X] \otimes_R R[Y]$（例~\ref{ex:poly-tensor}）.}：
\[
  r_1(m_1 \otimes n_1) + r_2(m_2 \otimes n_2) + \dots + r_k(m_k \otimes n_k),
\]
其中 $k \geq 1$，$r_i \in R$，$m_i \in M$ 且 $n_i \in N$.由于 $r_i(m_i \otimes n_i) = (r_i m_i) \otimes n_i$，我们可以将 $r_i m_i$ 重新记作 $m_i$，从而将上述线性组合改写为简单的和式：
\begin{equation}
  m_1 \otimes n_1 + m_2 \otimes n_2 + \dots + m_k \otimes n_k. \label{eq:1.4}
\end{equation}

在直和 $M \oplus N$ 中，相等性很容易定义：$(m, n) = (m', n')$ 当且仅当 $m = m'$ 且 $n = n'$.而在 $M \otimes_R N$ 中，形如~\eqref{eq:1.4} 的两个和式何时相等？在我们给出的张量积描述下，这并不容易判定，除了一种特殊情况：$M$ 与 $N$ 是自由 $R$-模，分别具有基 $\{e_i\}$ 与 $\{f_j\}$.在这种情况下，$M \otimes_R N$ 是以 $\{e_i \otimes f_j\}$ 为基的自由模，因此 $M \otimes_R N$ 中的每个元素都是（有限）和式 $\sum_{i,j} c_{ij} e_i \otimes f_j$（其中 $c_{ij} \in R$），且两个此类和式相等当且仅当同类项的系数相等.

为了在 $M$ 与 $N$ 没有基时描述 $M \otimes_R N$ 中的相等性，我们将利用 $M \otimes_R N$ 的\newterm{ump}{普遍映射性质}{universal mapping property}.张量积是代数学中第一个唯有通过普遍映射性质才能获得一致、合理理解的概念：$M \otimes_R N$ 是把 $M \times N$ 上的双线性映射转化为线性映射的万有对象.正如 Jeremy Kun \cite{kun} 所写，$M \otimes_R N$ 是所有从 $M \times N$ 出发的双线性映射的``守门人''（gatekeeper）.

在第~\ref{sec:bilinear}~节讨论双线性（与多重线性）映射之后，张量积的严格定义与构造将在第~\ref{sec:construction}~节中给出.张量积的具体实例见第~\ref{sec:examples}~节.在第~\ref{sec:properties}~节中，我们将展示张量积如何与模的其他构造相互作用.第~\ref{sec:base-ext}~节介绍\newterm{base-extension}{基变换}{base extension}（或称标量扩张）这一重要运算，它是利用张量积将 $R$-模转化为 $S$-模的过程，其中 $S$ 是另一个环.最后，在第~\ref{sec:physics}~节中，我们将介绍物理学中用于张量的指标记号.

以下是关于张量与张量积的简短历史.张量（tensor）一词源于拉丁语 \textit{tendere}，意为``拉伸''.1822 年，Cauchy 在连续介质力学中引入了柯西应力张量；1861 年，Riemann 在几何学中创立了黎曼曲率张量，但当时他们均未使用``张量''这一名称.1884 年，Gibbs \cite[第~3~章]{gibbs-vector} 引入了 $\R^3$ 中向量的张量积，称其为``不定积''（indeterminate product）\footnote{Gibbs 选用这一名称是因为用他的话来说，该乘积是``两个向量乘积的最一般形式''，因为它除了双线性之外不服从任何定律，而双线性是任何配称为乘积的运算所必须满足的.早在 1844 年，Grassmann 就创立了一种称为``开积''（open product）的特殊张量 \cite[第~3~章]{grassmann}.}，并将其应用于研究物体的应变.他在 1886 年将不定积推广到了 $n$ 维空间 \cite{gibbs-grassmann}.

Voigt 于 1898 年利用张量来描述晶体上的应力与应变 \cite{voigt}，``张量''一词在此处首次以其现代物理含义登场\footnote{用 $i, j, k$ 表示 $\R^3$ 的标准基，Gibbs 将形如 $a i \otimes i + b j \otimes j + c k \otimes k$（其中 $a, b, c > 0$）的和式称为正张量（right tensor）\cite[第~57~页]{gibbs-vector}，但我不知道这是否对 Voigt 的术语产生过影响.}.在几何学中，Ricci 在 19 世纪末使用了张量，他与 Levi-Civita 于 1901 年发表的论文 \cite{ricci-levicivita}（英译本见 \cite{ricci-english}）对 Einstein 创立广义相对论起到了关键作用.``张量''一词在物理学与数学中的广泛使用归功于 Einstein；而 Ricci 与 Levi-Civita 之前则用平淡的名称``体系''（systems）来称呼张量.Weyl 关于群论与量子力学的名著将向量空间 $V$ 与 $W$ 的张量积称为``乘积空间''（product space），并记作 $V \times W$ \cite[第~II~章第~10~节]{weyl-groups}.记号 $\otimes$ 归功于 Murray 与 von Neumann 在 1936 年关于 Hilbert 空间张量积（他们称之为``直积''）的研究 \cite[第~II~章]{murray-vonneumann}\footnote{感谢 Jim Casey 让我注意到文献 \cite{murray-vonneumann}.}.阿贝尔群 $A$ 与 $B$ 的张量积（使用了该名称，但记作 $A \circ B$ 而非 $A \otimes_\Z B$）由 Whitney 于 1938 年给出 \cite{whitney}.而交换环上模的张量积则由 Bourbaki 于 1948 年正式确立 \cite{bourbaki-alg}.

% sections/sec2_bilinear.tex
\section{双线性映射}
\label{sec:bilinear}

我们在前一节中已经将 $M \otimes_R N$ 的元素描述为服从规则~\eqref{eq:1.2} 与~\eqref{eq:1.3} 的和式~\eqref{eq:1.4}.其核心意图在于使得 $M \otimes_R N$ 成为满足规则~\eqref{eq:1.2} 与~\eqref{eq:1.3} 的``最自由''对象.而规则~\eqref{eq:1.2} 与~\eqref{eq:1.3} 的本质就是双线性（bilinearity）.这究竟是什么含义？

设 $M$、$N$ 与 $P$ 为 $R$-模.如果函数 $B \colon M \times N \to P$ 在固定另一个变元时对每个变元都是线性的（即 $R$-线性的），则称 $B$ 为\newterm{bilinear-map}{双线性映射}{bilinear map}（或双线性函数）：
\begin{align}
  B(m_1 + m_2, n) &= B(m_1, n) + B(m_2, n), & B(rm, n) &= rB(m, n), \label{eq:2.1} \\
  B(m, n_1 + n_2) &= B(m, n_1) + B(m, n_2), & B(m, rn) &= rB(m, n). \label{eq:2.2}
\end{align}
因此对每个 $n \in N$，$B(-, n)$ 是一个线性映射 $M \to P$；对每个 $m \in M$，$B(m, -)$ 是一个线性映射 $N \to P$.特别地，总有 $B(0, n) = 0$ 以及 $B(m, 0) = 0$.

以下是双线性映射的一些常见实例：
\begin{enumerate}
  \item $\R^n$ 上的点积 $v \cdot w$ 是双线性函数 $\R^n \times \R^n \to \R$.更一般地，对任意矩阵 $A \in \mathrm{M}_n(\R)$，函数 $\langle v, w \rangle = v \cdot Aw$ 都是双线性映射 $\R^n \times \R^n \to \R$.
  \item 矩阵乘法 $\mathrm{M}_{m,n}(\R) \times \mathrm{M}_{n,p}(\R) \to \mathrm{M}_{m,p}(\R)$ 是双线性的.点积是 $m = p = 1$ 的特例（此时将 $v \cdot w$ 写作 $v^\top w$）.
  \item 叉积 $v \times w$ 是双线性函数 $\R^3 \times \R^3 \to \R^3$.
  \item 行列式 $\det \colon \mathrm{M}_2(R) \to R$ 作为矩阵两列的函数是双线性的.
  \item 对于 $R$-模 $M$，标量乘法 $R \times M \to M$ 是双线性的.
  \item 环自身的乘法 $R \times R \to R$ 是双线性的.
  \item 设 $M$ 的对偶模为 $M^\vee = \Hom_R(M, R)$.由 $(\varphi, m) \mapsto \varphi(m)$ 给出的对偶配对 $M^\vee \times M \to R$ 是双线性的.
  \item 对于 $\varphi \in M^\vee$ 与 $\psi \in N^\vee$，映射 $M \times N \to R$（其中 $(m, n) \mapsto \varphi(m)\psi(n)$）是双线性的.
  \item 若 $M \times N \xrightarrow{B} P$ 是双线性的，且 $P \xrightarrow{L} Q$ 是线性的，则复合映射 $M \times N \xrightarrow{L \circ B} Q$ 也是双线性的.（这是一个至关重要的性质，请务必验证！）
  \item 根据第~\ref{sec:intro}~节的讨论，表达式 $m \otimes n$ 在 $m$ 与 $n$ 上应当是双线性的.换言之，我们希望由 $(m, n) \mapsto m \otimes n$ 给出的函数 $M \times N \to M \otimes_R N$ 是双线性的.
\end{enumerate}

以下是几个双元函数但\textbf{并非}双线性映射的典型反例：
\begin{enumerate}
  \item 对于 $R$-模 $M$，由 $(m, m') \mapsto m + m'$ 给出的加法运算 $M \times M \to M$ 通常不是双线性的：当 $m'$ 固定时它通常对 $m$ 不是加性的（即一般而言 $(m_1 + m_2) + m' \neq (m_1 + m') + (m_2 + m')$），当 $m$ 固定时它对 $m'$ 亦非加性.
  \item 对于 $\varphi \in M^\vee$ 与 $\psi \in N^\vee$，由 $(m, n) \mapsto \varphi(m) + \psi(n)$ 给出的和式 $M \times N \to R$ 通常不是双线性的.
  \item 将 $\mathrm{M}_n(\C)$ 视为 $\C$-向量空间.由 $(A, B) \mapsto A\bar{B}$ 给出的函数 $\mathrm{M}_n(\C) \times \mathrm{M}_n(\C) \to \mathrm{M}_n(\C)$ 不是双线性的.它是双加性的（即当固定另一个分量时对每个分量均可加），但观察标量乘法在第二个分量上的表现：对 $z \in \C$，$A\overline{zB} = \bar{z}(A\bar{B})$，而不是 $z(A\bar{B})$.
\end{enumerate}

为什么我们要把双线性函数写成 $B \colon M \times N \to P$，而不是写成 $B \colon M \oplus N \to P$？因为我们并没有在定义域上使用模的直和加法结构.这就是为什么对于加法 $R \times R \to R$ 与乘法 $R \times R \to R$，我们可以把加法写成 $R \oplus R \to R$（它在直和 $R \oplus R$ 上是线性的），但绝不能对乘法这样做：乘法是双线性的，却不是线性的（例如一般而言 $(r + r')(s + s') \neq rs + r's'$）.线性函数是广义的加法，而双线性函数是广义的乘法.切勿将 $M \times N$ 上的双线性函数与 $M \oplus N$ 上的线性函数混淆.

对于 $R$-模 $M_1, \dots, M_k$，如果函数 $f \colon M_1 \times \dots \times M_k \to M$ 在固定其他坐标时对每个 $m_i$ 都是线性的（即 $R$-线性的），则称 $f$ 为\newterm{multilinear-map}{多重线性映射}{multilinear map}（或 $k$-重线性映射）.因此 $2$-重线性即指双线性.以下是多重线性函数的一些例子：
\begin{enumerate}
  \item 标量三重积（混合积）$u \cdot (v \times w)$ 是三线性函数 $\R^3 \times \R^3 \times \R^3 \to \R$.
  \item 函数 $f(u, v, w) = (u \cdot v)w$ 是三线性函数 $\R^n \times \R^n \times \R^n \to \R^n$.
  \item 由 $(\varphi, m, n) \mapsto \varphi(m)n$ 给出的函数 $M^\vee \times M \times N \to N$ 是三线性的.
  \item 若 $B \colon M \times N \to P$ 与 $B' \colon P \times Q \to T$ 均为双线性的，则由 $(m, n, q) \mapsto B'(B(m, n), q)$ 给出的映射 $M \times N \times Q \to T$ 是三线性的.
  \item 具有 $k$ 个因子的乘法 $R \times \dots \times R \to R$ 是 $k$-重线性的.
  \item 行列式 $\det \colon \mathrm{M}_n(R) \to R$ 作为矩阵各列的函数是 $n$-重线性的.
  \item 若 $M_1 \times \dots \times M_k \xrightarrow{f} M$ 是 $k$-重线性的，且 $M \xrightarrow{L} N$ 是线性的，则复合映射 $M_1 \times \dots \times M_k \xrightarrow{L \circ f} N$ 也是 $k$-重线性的.
  \item 对于光滑映射 $f \colon \R^m \to \R^n$ 及点 $x \in \R^m$，$k$ 阶导数 $(D^k f)_x \colon (\R^m)^k \to \R^n$ 是 $k$-重线性的.（若不熟悉将一阶导数视作线性映射的观点，可忽略此例.）
\end{enumerate}

从 $M$ 到 $N$ 的所有 $R$-线性映射在函数的点态加法与标量数乘下构成一个 $R$-模 $\Hom_R(M, N)$.类似地，从 $M \times N$ 到 $P$ 的所有 $R$-双线性映射以相同方式构成一个 $R$-模 $\mathrm{Bil}_R(M, N; P)$.然而，与线性映射相比，双线性映射缺失了若干关键特征：
\begin{enumerate}
  \item 双线性映射 $M \times N \to P$ 不存在``核''（kernel）的概念，因为定义域 $M \times N$ 并非一个模（我们没有赋予其模运算使其保持线性）.
  \item 双线性映射 $M \times N \to P$ 的像集不一定构成子模.
\end{enumerate}

\begin{example}\label{ex:rank-one}
定义 $B \colon \R^n \times \R^n \to \mathrm{M}_n(\R)$ 为 $B(v, w) = v w^\top$，其中 $v$ 与 $w$ 均为列向量，因此 $v w^\top$ 是一个 $n \times n$ 矩阵.例如当 $n = 2$ 时，
\[
  B\left(\begin{pmatrix} a_1 \\ a_2 \end{pmatrix}, \begin{pmatrix} b_1 \\ b_2 \end{pmatrix}\right) = \begin{pmatrix} a_1 \\ a_2 \end{pmatrix} (b_1, b_2) = \begin{pmatrix} a_1 b_1 & a_1 b_2 \\ a_2 b_1 & a_2 b_2 \end{pmatrix}.
\]
一般地，若用 $\R^n$ 的标准基表示 $v = \sum a_i e_i$ 与 $w = \sum b_j e_j$，则 $v w^\top$ 就是 $n \times n$ 矩阵 $(a_i b_j)$.该矩阵关于 $v$ 和 $w$ 分别都是 $\R$-线性的，所以 $B$ 是双线性的.

然而当 $n \geq 2$ 时，$B$ 的像在加法下并不封闭，因此它不是 $\mathrm{M}_n(\R)$ 的子空间.原因何在？由于矩阵 $B(v, w)$ 的每一列都是列向量 $v$ 的标量倍数，所以其秩至多为 $1$（非零时秩恰为 $1$，零时秩为 $0$）.考虑矩阵
\[
  B(e_1, e_1) + B(e_2, e_2) = e_1 e_1^\top + e_2 e_2^\top = \begin{pmatrix}
    1 & 0 & \dots & 0 \\
    0 & 1 & \dots & 0 \\
    \vdots & \vdots & \ddots & \vdots \\
    0 & 0 & \dots & 0
  \end{pmatrix},
\]
该矩阵的像空间是 $2$ 维的（秩为 $2$），因此对于 $\R^n$ 中的任何向量 $v$ 与 $w$，都不可能有 $B(e_1, e_1) + B(e_2, e_2) = B(v, w)$.（类似地，$\sum_{i=1}^n B(e_i, e_i)$ 是 $n \times n$ 单位矩阵，它绝不能写成单个矩阵 $B(v, w)$ 的形式.）
\end{example}

% sections/sec3_construction.tex
\section{张量积的构造}
\label{sec:construction}

到达 $R$-模 $P$ 的双线性映射 $M \times N \to P$ 可以与线性映射 $P \to Q$ 复合，从而得到一个双线性映射 $M \times N \to Q$：
\begin{center}
\begin{tikzcd}[row sep=1.8cm, column sep=2.5cm]
M \times N \arrow[r, "B\text{（双线性）}"] \arrow[dr, "L \circ B\text{（复合仍为双线性）}"'] & P \arrow[d, "L\text{（线性）}"] \\
& Q
\end{tikzcd}
\end{center}

我们将把 $M$ 与 $N$ 的张量积构造为一个普遍映射问题的解：寻找一个 $R$-模 $T$ 与一个双线性映射 $b \colon M \times N \to T$，使得 $M \times N$ 上的每一个双线性映射都可以唯一地分解为该双线性映射 $b$ 与一个从 $T$ 出发的线性映射的复合：
\begin{center}
\begin{tikzcd}[row sep=1.8cm, column sep=2.5cm]
M \times N \arrow[r, "b\text{（双线性）}"] \arrow[dr, "B\text{（双线性）}"'] & T \arrow[d, "\exists\,L\text{（线性？）}", dashed] \\
& P
\end{tikzcd}
\end{center}

这完全类似于群 $G$ 的交换化 $G/[G, G]$ 的普遍映射性质：从 $G$ 到阿贝尔群 $A$ 的同态与从 $G/[G, G]$ 到 $A$ 的同态是``同一回事''，因为每个同态 $f \colon G \to A$ 都可以唯一地分解为规范同态 $\pi \colon G \to G/[G, G]$ 与唯一同态 $\widetilde{f} \colon G/[G, G] \to A$ 的复合：
\begin{center}
\begin{tikzcd}[row sep=1.8cm, column sep=2.5cm]
G \arrow[r, "\pi"] \arrow[dr, "f"'] & G/[G, G] \arrow[d, "\exists!\,\widetilde{f}", dashed] \\
& A
\end{tikzcd}
\end{center}

\begin{definition}\label{def:tensor-product}
$M$ 与 $N$ 的\newterm{tensor-prod-def}{张量积}{tensor product} $M \otimes_R N$ 是一个赋予了双线性映射 $M \times N \xrightarrow{\otimes} M \otimes_R N$ 的 $R$-模，满足：对于每一个双线性映射 $M \times N \xrightarrow{B} P$，都存在唯一的线性映射 $M \otimes_R N \xrightarrow{L} P$ 使得如下交换图成立：
\begin{equation}
\begin{tikzcd}[row sep=large, column sep=large]
M \times N \arrow[r, "\otimes"] \arrow[dr, "B"'] & M \otimes_R N \arrow[d, "\exists!\,L", dashed] \\
& P
\end{tikzcd} \label{eq:3.0}
\end{equation}
即对所有 $m \in M$ 与 $n \in N$ 均有 $B(m, n) = L(m \otimes n)$.
\end{definition}

尽管群交换化 $G/[G, G]$ 普遍映射性质中的函数全都是群同态（分别从 $G$ 与 $G/[G, G]$ 出发），但在 $M \otimes_R N$ 的普遍映射性质中，所涉及的函数类型并不相同：从 $M \times N$ 出发的函数是\textbf{双线性}的，而从 $M \otimes_R N$ 出发的函数是\textbf{线性}的.简而言之，从 $M \times N$ 出发的双线性映射被转化为了从 $M \otimes_R N$ 出发的线性映射.

张量积的定义不仅包含一个新模 $M \otimes_R N$，还包含一个到达它的特指双线性映射 $\otimes \colon M \times N \to M \otimes_R N$.这也类似于群交换化 $G/[G, G]$ 的普遍映射性质：不仅需要群 $G/[G, G]$，还需要固定同态 $\pi \colon G \to G/[G, G]$，使得所有到达阿贝尔群的同态都能通过它分解.普遍映射性质必须固定这笔额外的结构信息.

在具体构造张量积之前，我们先证明任意两个张量积在本质上都是相同的.设 $R$-模 $T$ 与 $T'$，以及双线性映射 $M \times N \xrightarrow{b} T$ 与 $M \times N \xrightarrow{b'} T'$，均满足张量积的普遍映射性质.由 $M \times N \xrightarrow{b} T$ 的万有性，映射 $M \times N \xrightarrow{b'} T'$ 唯一地通过 $T$ 分解：存在唯一的线性映射 $f \colon T \to T'$ 使得
\begin{equation}
\begin{tikzcd}[row sep=large, column sep=large]
M \times N \arrow[r, "b"] \arrow[dr, "b'"'] & T \arrow[d, "f"] \\
& T'
\end{tikzcd} \label{eq:3.1}
\end{equation}
交换.反过来，由 $M \times N \xrightarrow{b'} T'$ 的万有性，映射 $M \times N \xrightarrow{b} T$ 也唯一地通过 $T'$ 分解：存在唯一的线性映射 $f' \colon T' \to T$ 使得
\begin{equation}
\begin{tikzcd}[row sep=large, column sep=large]
M \times N \arrow[r, "b'"] \arrow[dr, "b"'] & T' \arrow[d, "f'"] \\
& T
\end{tikzcd} \label{eq:3.2}
\end{equation}
交换.将~\eqref{eq:3.1} 与~\eqref{eq:3.2} 串联拼接成如下交换图：
\begin{center}
\begin{tikzcd}[row sep=large, column sep=large]
& T \arrow[d, "f"] \\
M \times N \arrow[ur, "b"] \arrow[r, "b'"{description}] \arrow[dr, "b"'] & T' \arrow[d, "f'"] \\
& T
\end{tikzcd}
\end{center}
抹去中间层，便得到交换图：
\begin{equation}
\begin{tikzcd}[row sep=large, column sep=large]
M \times N \arrow[r, "b"] \arrow[dr, "b"'] & T \arrow[d, "f' \circ f"] \\
& T
\end{tikzcd} \label{eq:3.3}
\end{equation}
根据 $(T, b)$ 的万有性，使得图~\eqref{eq:3.3} 交换的从 $T$ 到 $T$ 的线性映射是唯一的.而恒等映射 $\id_T$ 显然满足条件，故必有 $f' \circ f = \id_T$.同理，按相反顺序将~\eqref{eq:3.1} 与~\eqref{eq:3.2} 串联拼接，便可得到 $f \circ f' = \id_{T'}$.因此，$T$ 与 $T'$ 通过 $f$ 是同构的 $R$-模，且 $f \circ b = b'$，这意味着 $f$ 将 $b$ 与 $b'$ 相结合.因此，$M$ 与 $N$ 的任意两个张量积都可以以唯一的方式相识同构，并且这种同构与从 $M \times N$ 出发的特指双线性映射是相容的\footnote{``相容''的意思是，若 $f \colon T \to T'$ 是使 $f \circ b = b'$ 的唯一同构，则对任意 $(m, n)$ 均有 $f(b(m, n)) = b'(m, n)$.}.

\begin{theorem}\label{thm:tensor-exists}
$M$ 与 $N$ 的张量积必定存在.
\end{theorem}

\begin{proof}
将笛卡尔积 $M \times N$ 单纯视为一个集合.我们在该集合上构造自由 $R$-模：
\[
  F_R(M \times N) = \bigoplus_{(m, n) \in M \times N} R \delta_{(m, n)}.
\]
（这是一个极其庞大的 $R$-模.若 $R = \R$ 且 $M = N = \R^3$，则 $F_R(M \times N)$ 是由 $\R^6$ 多个 $R$ 的副本构成的直和.这个直和遍历 $\R^3$ 中所有成对的向量，绝非仅仅取自一组基，其分量位于 $R$ 中.对于一般的模，基甚至根本不存在.）

设 $D$ 是 $F_R(M \times N)$ 中由以下所有形如这些关系的元素所生成的子模：
\begin{gather*}
  \delta_{(m + m', n)} - \delta_{(m, n)} - \delta_{(m', n)}, \quad
  \delta_{(m, n + n')} - \delta_{(m, n)} - \delta_{(m, n')}, \\
  \delta_{(rm, n)} - r\delta_{(m, n)}, \quad
  \delta_{(m, rn)} - r\delta_{(m, n)}.
\end{gather*}
模 $F_R(M \times N)$ 对 $D$ 的商模便可用作张量积：令
\[
  M \otimes_R N := F_R(M \times N) / D.
\]
我们将 $M \otimes_R N$ 中的陪集 $\delta_{(m, n)} + D$ 记作 $m \otimes n$.

根据 $D$ 的定义，在商模 $F_R(M \times N) / D$ 中有同余关系：
\[
  \delta_{(m + m', n)} \equiv \delta_{(m, n)} + \delta_{(m', n)} \pmod D,
\]
在 $M \otimes_R N$ 中这恰好写成
\[
  (m + m') \otimes n = m \otimes n + m' \otimes n.
\]
同理，在 $M \otimes_R N$ 中亦有 $m \otimes (n + n') = m \otimes n + m \otimes n'$ 以及 $r(m \otimes n) = rm \otimes n = m \otimes rn$.这些关系正是当初定义子模 $D$ 的原因，它们说明由 $(m, n) \mapsto m \otimes n$ 给出的函数 $M \times N \xrightarrow{\otimes} M \otimes_R N$ 确为双线性映射.（除此映射之外，我们不会考虑任何其他从 $M \times N$ 到 $M \otimes_R N$ 的函数.）

现在我们证明从 $M \times N$ 出发的所有双线性映射都可以唯一地通过上面写出的双线性映射 $M \times N \to M \otimes_R N$ 分解.假设 $P$ 是一个 $R$-模，且 $M \times N \xrightarrow{B} P$ 是一双线性映射.若单纯将 $M \times N$ 视为一个集合（暂时忽略 $B$ 的双线性），由自由模的普遍映射性质，函数 $B \colon M \times N \to P$ 可以唯一延拓为一个线性映射 $\ell \colon F_R(M \times N) \to P$，满足 $\ell(\delta_{(m, n)}) = B(m, n)$，使得交换图成立：
\begin{center}
\begin{tikzcd}[row sep=large, column sep=large]
M \times N \arrow[r, "\delta"] \arrow[dr, "B"'] & F_R(M \times N) \arrow[d, "\ell"] \\
& P
\end{tikzcd}
\end{center}
我们希望说明 $\ell$ 在商模 $M \otimes_R N$ 上能够良定义，这等价于证明 $\ker \ell$ 包含子模 $D$.由 $B$ 的双线性：
\begin{gather*}
  B(m + m', n) = B(m, n) + B(m', n), \quad B(m, n + n') = B(m, n) + B(m, n'), \\
  rB(m, n) = B(rm, n) = B(m, rn),
\end{gather*}
代入 $\ell$ 的定义即得：
\begin{gather*}
  \ell(\delta_{(m + m', n)}) = \ell(\delta_{(m, n)}) + \ell(\delta_{(m', n)}), \quad
  \ell(\delta_{(m, n + n')}) = \ell(\delta_{(m, n)}) + \ell(\delta_{(m, n')}), \\
  r\ell(\delta_{(m, n)}) = \ell(\delta_{(rm, n)}) = \ell(\delta_{(m, rn)}).
\end{gather*}
由于 $\ell$ 是线性的，这些等式说明 $\ell$ 将子模 $D$ 的每一组生成元都映为 $0$.因此 $\ker \ell$ 包含子模 $D$，从而 $\ell$ 诱导出唯一的商模线性映射 $L \colon F_R(M \times N)/D \to P$，使得 $L(\delta_{(m, n)} + D) = \ell(\delta_{(m, n)}) = B(m, n)$.

由于 $F_R(M \times N)/D = M \otimes_R N$ 且 $\delta_{(m, n)} + D = m \otimes n$，这正是交换图：
\begin{equation}
\begin{tikzcd}[row sep=large, column sep=large]
M \times N \arrow[r, "\otimes"] \arrow[dr, "B"'] & M \otimes_R N \arrow[d, "L"] \\
& P
\end{tikzcd} \label{eq:3.4}
\end{equation}
这表明从 $M \times N$ 出发的任何双线性映射 $B$ 均来自于一个线性映射 $L \colon M \otimes_R N \to P$，使得对所有 $m \in M$ 与 $n \in N$ 均有 $L(m \otimes n) = B(m, n)$.

最后需证明满足图~\eqref{eq:3.4} 交换的线性映射 $L$ 是唯一的.回到 $M \otimes_R N$ 作为自由模商模的定义，自由模 $F_R(M \times N)$ 中的每个元素都是形如 $\sum_{i=1}^k r_i \delta_{(m_i, n_i)}$ 的有限和.商映射 $F_R(M \times N) \to F_R(M \times N)/D = M \otimes_R N$ 是线性的，因此 $M \otimes_R N$ 中的每一个元素都是有限和：
\begin{equation}
  r_1 (m_1 \otimes n_1) + \dots + r_k (m_k \otimes n_k). \label{eq:3.5}
\end{equation}
这意味着元素 $m \otimes n$ 在 $R$ 上张成了整个 $M \otimes_R N$.因此，从 $M \otimes_R N$ 出发的任何线性映射由其在所有元素 $m \otimes n$ 上的取值完全且唯一决定.故而满足 $m \otimes n \mapsto B(m, n)$ 的线性映射 $M \otimes_R N \to P$ 至多只有一个.我们在图~\eqref{eq:3.4} 中已经构造出了这样一个线性映射，因此它必定是唯一的.
\end{proof}

既然证明了 $M$ 与 $N$ 的张量积存在\footnote{1948 年，Bourbaki \cite{bourbaki-alg} 像我们这里一样，将张量积构造为自由模的商模.1938 年，Whitney \cite{whitney} 对阿贝尔群的张量积给出了等价构造.}，其本质唯一性使得我们称 $M \otimes_R N$ 为``\textbf{那个}''张量积（the tensor product），而非``一个''张量积.不要忘记，这一构造不仅包含模 $M \otimes_R N$，还包含由 $(m, n) \mapsto m \otimes n$ 给出的特指双线性映射 $M \times N \xrightarrow{\otimes} M \otimes_R N$，从 $M \times N$ 出发的所有双线性映射都必须通过它分解.我们称这个特指映射为从 $M \times N$ 到张量积的规范双线性映射.

$M \otimes_R N$ 中的元素称为\termcn{张量}，通常用字母 $t$ 表示.具有形如 $m \otimes n$ 的张量称为\newterm{elementary-tensor}{简单张量}{elementary tensor}（也称初等张量、纯张量、单项张量或可分解张量）.正如集合 $A$ 上的自由模 $F_R(A)$ 中的元素通常不是单个 $\delta_a$ 而是它们的线性组合一样，$M \otimes_R N$ 中的元素通常不是简单张量\footnote{对于 $R \neq 0$，$R^2 \otimes_R R^2$ 中非简单张量的一个显式例子将在例~\ref{ex:nonelementary} 中给出.我们实际上在例~\ref{ex:rank-one} 中就已经见过一个：$e_1 e_1^\top + e_2 e_2^\top$ 无法写成单个 $v w^\top$ 的形式.}，而是简单张量的线性组合.事实上，由于 $r(m \otimes n) = (rm) \otimes n$，每一个张量都可以写成若干简单张量的和式.这表明 $M \otimes_R N$ 中的所有元素均具有式~\eqref{eq:1.4} 的形式.

每个张量都是简单张量的和，但其自身未必是简单张量，这一特性常令初学者感到困惑.困惑的一个根源在于：在直和 $M \oplus N$ 中每个元素都是一对 $(m, n)$，为什么在 $M \otimes_R N$ 中不能每个元素都具有 $m \otimes n$ 的形式呢？这里有两个相关的直观类比有助于消除这种陌生感：
\begin{itemize}
  \item $R$-模 $R[X, Y]$ 是 $R[X]$ 与 $R[Y]$ 的张量积（参见例~\ref{ex:poly-tensor}）.正如 Eisenbud 与 Harris 在其概型论著作中所指出的 \cite[第~39~页]{eisenbud-harris}，研究双变元多项式绝不仅限于研究形如 $f(X)g(Y)$ 的乘积多项式.事实上，$R[X, Y]$ 中的绝大多数多项式都不能写成 $f(X)g(Y)$，但它们全部是此类乘积的有限和（实际上是单项式 $a_{ij} X^i Y^j$ 的和）.
  \item 简单张量在全体张量中的地位，完全类似于二元偏微分方程（PDE）的分离变量解 $f(x)g(y)$ 在所有通解中的地位\footnote{感谢 David Vogan 提出这一精妙类比.}.PDE 的解通常是不可分离的；我们首先求出所有分离变量解，然后将一般通解表示为分离变量解的叠加和（或级数和）.
\end{itemize}

从现在起，请忘掉将 $M \otimes_R N$ 显式构造为庞大自由模 $F_R(M \times N)$ 之商的做法.试图通过商模的构造细节来思考张量积，只会带来不必要的混乱.更重要的是铭记张量积的普遍映射性质，我们从下一节起将系统地运用这一性质.为了熟悉 $\otimes$ 的双线性，我们先证明两个简单的命题：

\begin{theorem}\label{thm:spanning-tensor}
设 $M$ 与 $N$ 分别具有生成集 $\{x_i\}_{i \in I}$ 与 $\{y_j\}_{j \in J}$.则张量积 $M \otimes_R N$ 作为 $R$-模由所有简单张量 $x_i \otimes y_j$ 线性生成.
\end{theorem}

\begin{proof}
$M \otimes_R N$ 中的简单张量形如 $m \otimes n$.将它们展开为 $m = \sum_i a_i x_i$ 与 $n = \sum_j b_j y_j$（其中非零系数仅有有限个）.根据 $\otimes$ 的双线性，
\[
  m \otimes n = \left(\sum_i a_i x_i\right) \otimes \left(\sum_j b_j y_j\right) = \sum_{i, j} a_i b_j (x_i \otimes y_j)
\]
是诸张量 $x_i \otimes y_j$ 的线性组合.由于每个简单张量都是这组特定张量 $x_i \otimes y_j$ 的线性组合，而每个一般张量又是有限个简单张量的和，因此诸 $x_i \otimes y_j$ 构成了 $M \otimes_R N$ 的一个生成集.
\end{proof}

\begin{example}\label{ex:rk-tensor}
设 $e_1, \dots, e_k$ 为 $R^k$ 的标准基.则 $R$-模 $R^k \otimes_R R^k$ 由 $k^2$ 个简单张量 $e_i \otimes e_j$ 线性生成.我们在后文（定理~\ref{thm:free-tensor-basis}）将会看到，这些简单张量实际上构成了 $R^k \otimes_R R^k$ 的一组自由基；当 $R$ 为域时，这与第~\ref{sec:intro}~节中物理学家使用基对向量空间张量积所作的``定义''完全吻合.
\end{example}

\begin{theorem}\label{thm:zero-tensor}
在 $M \otimes_R N$ 中，恒有 $m \otimes 0 = 0$ 以及 $0 \otimes n = 0$.
\end{theorem}

\begin{proof}
这与环中证明 $a \cdot 0 = 0$ 完全相同：由于固定 $m$ 时 $m \otimes n$ 对 $n$ 是加性的，
\[
  m \otimes 0 = m \otimes (0 + 0) = m \otimes 0 + m \otimes 0.
\]
两边减去 $m \otimes 0$，即得 $m \otimes 0 = 0$.同理可证 $0 \otimes n = 0$.
\end{proof}

\begin{example}\label{ex:torsion-collapse}
若 $A$ 是有限阿贝尔群，则 $\Q \otimes_\Z A = 0$.这是因为其中的每个简单张量都为 $0$：对任意 $a \in A$，存在正整数 $n$ 使得 $na = 0$.于是在 $\Q \otimes_\Z A$ 中，对任意有理数 $r$，
\[
  r \otimes a = \left(n \cdot \frac{r}{n}\right) \otimes a = \frac{r}{n} \otimes (na) = \frac{r}{n} \otimes 0 = 0.
\]
每个张量都是简单张量的和，既然每个简单张量都为 $0$，那么全体张量皆为 $0$.（例如在 $\Q \otimes_\Z \Z/7\Z$ 中有 $\frac{1}{3} \otimes (5 \bmod 7) = 0$.由此可见，完全可能出现 $m \otimes n = 0$ 而 $m$ 与 $n$ 均非零的情形.）

要证明 $\Q \otimes_\Z A = 0$，甚至不需要 $A$ 是有限群，只需 $A$ 的每个元素都是阶有限的扭元即可.商群 $\Q/\Z$ 具有这一性质，因此 $\Q \otimes_\Z (\Q/\Z) = 0$.同理亦有 $\Q/\Z \otimes_\Z \Q/\Z = 0$.
\end{example}

由于 $M \otimes_R N$ 在加法上由简单张量生成，从 $M \otimes_R N$ 出发的任何线性（乃至纯加性）映射由它在简单张量上的取值完全决定.这正是为什么在实践中，张量积上的线性映射通常仅通过其在简单张量上的作用来指定.这好比用矩阵来描述有限维自由模之间的线性映射：矩阵表面上只给出了基向量的像，但这些信息已足以完全确定映射在整个空间上的行为.

然而，基向量与简单张量之间存在一个至关重要的本质区别：\textbf{简单张量之间存在大量的线性关系}.从 $\R^2$ 出发的线性映射由它在 $(1, 0), (2, 3), (8, 4)$ 与 $(-1, 5)$ 上的取值决定，但这四个取值绝不能任意给出，它们必须满足这四个向量之间已有的所有线性关系，因为线性映射必须保持线性关系.同样地，在简单张量上随意定义一个函数，通常\textbf{根本无法}延拓为张量积上的线性映射：简单张量虽然生成了张量积，但它们绝不是线性无关的！

不能随心所欲地通过二元函数在简单张量上构造映射.例如，令 $m \otimes m' \mapsto m + m'$ 或 $m \otimes m' \mapsto m$ 的所谓``映射'' $M \otimes_R M \to M$ 是完全不成立的，因为 $m \otimes m' = (-m) \otimes (-m')$，但一般而言 $m + m' \neq -m - m'$ 且 $m \neq -m$.从 $M \otimes_R N$ 出发构造线性映射的唯一正确途径，是利用 $M \otimes_R N$ 的普遍映射性质（将双线性映射转化为线性映射），因为简单张量之间所有可能的线性关系——无论显而易见的还是极其隐晦的——都已经内建在普遍映射性质之中了.我们在第~\ref{sec:examples}~节中将进行大量这方面的练习.理解如何利用普遍映射性质来计算具体例题和证明张量积的性质，是掌握张量积的唯一捷径；如果无法在 $M \otimes_R N$ 上正确构造函数，就谈不上真正理解了张量积.

张量积可以自然推广到多个因子的情形.给定 $k$ 个模 $M_1, \dots, M_k$，存在一个模 $M_1 \otimes_R \dots \otimes_R M_k$，它是 $k$-重线性映射的万有对象：它配有一个 $k$-重线性映射 $M_1 \times \dots \times M_k \xrightarrow{\otimes} M_1 \otimes_R \dots \otimes_R M_k$，且从 $M_1 \times \dots \times M_k$ 出发的任何 $k$-重线性映射都能唯一地通过与从 $M_1 \otimes_R \dots \otimes_R M_k$ 出发的线性映射复合而分解：
\begin{center}
\begin{tikzcd}[row sep=large, column sep=large]
M_1 \times \dots \times M_k \arrow[r, "\otimes"] \arrow[dr, "f"'] & M_1 \otimes_R \dots \otimes_R M_k \arrow[d, "\exists!\,L", dashed] \\
& P
\end{tikzcd}
\end{center}
$(m_1, \dots, m_k)$ 在该模中的像记作 $m_1 \otimes \dots \otimes m_k$.该 $k$-重张量积可以直接通过自由模 $F_R(M_1 \times \dots \times M_k)$ 的商模来构造，也可以通过每次对两个模作张量积来逐步构造：
\[
  (\dots((M_1 \otimes_R M_2) \otimes_R M_3) \dots) \otimes_R M_k.
\]
两种构造满足相同的普遍映射性质，因而在规范同构意义下是等价的.单个模 $M$ 的张量幂 $M^{\otimes k}$ 是多重张量积的重要特例：
\[
  M^{\otimes 0} = R, \quad M^{\otimes 1} = M, \quad M^{\otimes 2} = M \otimes_R M, \quad M^{\otimes 3} = M \otimes_R M \otimes_R M, \dots
\]
（其中约定 $M^{\otimes 0} = R$，正如约定 $a^0 = 1$ 一样.）

\subsection*{常见初学者问答}

现在我们解答关于张量积的七个初学者常见疑问：

\begin{enumerate}
  \item \textbf{究竟什么是 $m \otimes n$？}
  
  严格来说，$m \otimes n$ 就是有序对 $(m, n) \in M \times N$ 在张量积定义中的规范双线性映射 $M \times N \xrightarrow{\otimes} M \otimes_R N$ 之下的像.
  
  另一个在实践中更为切合的回答是：$m \otimes n$ 是张量积 $M \otimes_R N$ 中的这样一个元素——对于任意双线性映射 $M \times N \xrightarrow{B} P$，与其对应的线性映射 $M \otimes_R N \to P$ 在该元素处的值恰好等于 $B(m, n)$.请回顾定理~\ref{thm:tensor-exists} 的证明，确认这一性质的确成立.
  
  \item \textbf{$m \otimes n = 0$ 意味着什么？}
  
  $m \otimes n = 0$ 当且仅当从 $M \times N$ 出发的\textbf{每一个}双线性映射在 $(m, n)$ 处均取值为 $0$.
  
  事实上，若 $m \otimes n = 0$，则对任意双线性映射 $B \colon M \times N \to P$，存在与之相伴的线性映射 $L \colon M \otimes_R N \to P$，使得
  \[
    B(m, n) = L(m \otimes n) = L(0) = 0.
  \]
  反之，如果从 $M \times N$ 出发的每个双线性映射都把 $(m, n)$ 映为 $0$，那么规范双线性映射 $M \times N \xrightarrow{\otimes} M \otimes_R N$ 作为其中一个特例，必然也将 $(m, n)$ 映为 $0$.既然此映射把 $(m, n)$ 映为 $m \otimes n$，我们便断定 $m \otimes n = 0$.
  
  这一事实带来了一个极其重要的实用技巧：\textbf{要证明某个简单张量 $m \otimes n$ 不为零，只需找到一个双线性映射 $B$ 使得 $B(m, n) \neq 0$ 即可}.请牢记这一思想！它将在定理~\ref{thm:free-tensor-basis} 的证明中发挥关键作用.
  
  这也给出了定理~\ref{thm:zero-tensor}（$m \otimes 0 = 0$ 且 $0 \otimes n = 0$）的另一种证明，因为任意双线性映射显然满足 $B(m, 0) = 0$ 与 $B(0, n) = 0$.
  
  作为练习，请利用普遍映射性质证明：在 $M_1 \otimes_R \dots \otimes_R M_k$ 中 $m_1 \otimes \dots \otimes m_k = 0$ 当且仅当所有从 $M_1 \times \dots \times M_k$ 出发的 $k$-重线性映射都在 $(m_1, \dots, m_k)$ 处取值为 $0$.
  
  \item \textbf{$M \otimes_R N = 0$ 意味着什么？}
  
  张量积 $M \otimes_R N = 0$ 当且仅当从 $M \times N$ 出发的每一个双线性映射（到达任意模）都恒等于零.
  
  首先假设 $M \otimes_R N = 0$.则所有简单张量 $m \otimes n$ 均为 $0$，由第 (2) 问的结论，对所有双线性映射均有 $B(m, n) = 0$，从而 $B$ 恒为零.反之，若每一个双线性映射均恒为零，则规范双线性映射 $M \times N \xrightarrow{\otimes} M \otimes_R N$ 恒为零，这意味着对所有 $m$ 和 $n$ 均有 $m \otimes n = 0$.由于简单张量生成整个张量积，所有简单张量的消亡蕴含了整个模 $M \otimes_R N = 0$.
  
  回到例~\ref{ex:torsion-collapse}，当 $A$ 为扭群时 $\Q \otimes_\Z A = 0$，这无非是说从 $\Q \times A$ 出发的每一个 $\Z$-双线性映射都恒为零.这很容易直接验证：若 $B$ 是这样一个映射且 $na = 0$（$n \geq 1$），则
  \[
    B(r, a) = B\left(n \cdot \frac{r}{n}, a\right) = B\left(\frac{r}{n}, na\right) = B\left(\frac{r}{n}, 0\right) = 0.
  \]
  反用这一结论：\textbf{要证明某个张量积模 $M \otimes_R N \neq 0$，只需在 $M \times N$ 上构造出一个非零的双线性映射即可}.
  
  \item \textbf{$\sum_{i=1}^k m_i \otimes n_i = \sum_{j=1}^\ell m'_j \otimes n'_j$ 意味着什么？}
  
  $\sum_{i=1}^k m_i \otimes n_i = \sum_{j=1}^\ell m'_j \otimes n'_j$ 当且仅当对从 $M \times N$ 出发的所有双线性映射 $B$，都有
  \[
    \sum_{i=1}^k B(m_i, n_i) = \sum_{j=1}^\ell B(m'_j, n'_j).
  \]
  证明思路与前两问完全相同，留给读者作为练习.例如，$\sum_{i=1}^k m_i \otimes n_i = 0$ 意味着对所有双线性映射 $B$ 均有 $\sum_{i=1}^k B(m_i, n_i) = 0$.
  
  \item \textbf{张量积是如何让双线性映射看起来像线性映射的？}
  
  对于双线性映射 $B \colon M \times N \to P$，其双线性由式~\eqref{eq:2.1} 与~\eqref{eq:2.2} 刻画.对于与之对应的线性映射 $L \colon M \otimes_R N \to P$，由于
  \begin{gather*}
    (m_1 + m_2) \otimes n = m_1 \otimes n + m_2 \otimes n, \quad m \otimes (n_1 + n_2) = m \otimes n_1 + m \otimes n_2, \\
    (rm) \otimes n = r(m \otimes n) = m \otimes (rn),
  \end{gather*}
  在 $L$ 上这四个关于 $B$ 的双线性等式全部变成了 $L(t + t') = L(t) + L(t')$ 与 $L(rt) = rL(t)$ 的特例，这正是 $L$ 为线性映射的定义.
  
  \item \textbf{张量积在数学之外有哪些应用？}
  
  张量广泛应用于物理学与工程学（连续介质应力、弹性力学、电磁学、广义相对论度规、弥散磁共振成像等），在这些领域中，张量是在坐标变换下服从多重线性变换规则的物理量.物理学中的张量处理将在第~\ref{sec:physics}~节中详细讨论.
  
  \item \textbf{有没有什么办法能在头脑中直观想象张量积？}
  
  张量（哪怕是简单张量）并没有像``向量是一根带箭头的线段''那样直观的几何图像.上一问提到的物理表现虽然具体，但对理解抽象模的张量积并无直接帮助.
  
  初学者在面对张量积时都会感到不适应.Cathy O'Neil 与 Johan de Jong 的两段名言精彩地概括了理解张量积的心路历程\footnote{分别参见 \url{https://mathbabe.org/2011/07/20/tensor-product/} 与 MathOverflow 讨论 \url{https://mathoverflow.net/questions/35466}.}：
  \begin{itemize}
    \item \textbf{O'Neil}：``几个月后，我忽然意识到了一件事.我其实并没有比以前更'理解'张量积，但我已经'习惯'于不去理解它们了.这真是太奇妙了.每当张量积出现时，我不再感到痛苦焦虑；相反，我几乎被它们精巧狡黠的运作方式逗乐了.''
    \item \textbf{de Jong}：``正是那些你能证明的具体定理，塑造了你应当如何思考张量积.换言之，你要让初等的引理和具体的例子来逐步构建你对该数学对象的直觉.除此之外别无他法，根本不会有某种神奇的直觉突然从天而降来帮你'秒懂'它.''
  \end{itemize}
\end{enumerate}

\begin{remark}\label{rem:whitney}
最早在向量空间之外定义张量积的 Hassler Whitney，在阿贝尔群 $A$ 与 $B$ 存在到达阿贝尔群 $C$ 的 $\Z$-双线性映射 $A \times B \to C$ 时，称它们为关于 $C$ 的群对（group pair），并写道 \cite[第~499~页]{whitney}：``任何这样的群对都可以通过选取一个同态 $A \otimes_\Z B \to C$ 来定义''.因此，把 $\otimes_\Z$ 视作解决普遍映射问题的思想，本质上正是归功于 Whitney.
\end{remark}

% sections/sec4_examples.tex
\section{张量积的实例}
\label{sec:examples}

\begin{theorem}\label{thm:cyclic-tensor}
设 $a$ 与 $b$ 为正整数，令 $d = \gcd(a, b)$.则作为阿贝尔群有
\[
  \Z/a\Z \otimes_\Z \Z/b\Z \cong \Z/d\Z.
\]
特别地，$\Z/a\Z \otimes_\Z \Z/b\Z = 0$ 当且仅当 $\gcd(a, b) = 1$.
\end{theorem}

\begin{proof}
由于 $1$ 生成 $\Z/a\Z$ 与 $\Z/b\Z$，根据定理~\ref{thm:spanning-tensor}，$1 \otimes 1$ 生成 $\Z/a\Z \otimes_\Z \Z/b\Z$.又由
\[
  a(1 \otimes 1) = a \otimes 1 = 0 \otimes 1 = 0 \quad\text{且}\quad b(1 \otimes 1) = 1 \otimes b = 1 \otimes 0 = 0,
\]
可知 $1 \otimes 1$ 的加法阶整除 $a$ 且整除 $b$，从而整除 $d = \gcd(a, b)$.因此 $|\Z/a\Z \otimes_\Z \Z/b\Z| \leq d$.

为了证明 $\Z/a\Z \otimes_\Z \Z/b\Z$ 的阶至少为 $d$，我们在 $\Z/a\Z \otimes_\Z \Z/b\Z$ 上构造一个满的 $\Z$-线性映射到达 $\Z/d\Z$.由于 $d \mid a$ 且 $d \mid b$，自然存在典范约化 $\Z/a\Z \to \Z/d\Z$ 与 $\Z/b\Z \to \Z/d\Z$.考虑映射
\[
  B \colon \Z/a\Z \times \Z/b\Z \to \Z/d\Z, \quad B(x \bmod a, y \bmod b) = xy \bmod d.
\]
这在每个因子上先取模 $d$ 再相乘，它是显然良定义且 $\Z$-双线性的.因此由张量积的普遍映射性质，存在唯一的 $\Z$-线性映射 $f \colon \Z/a\Z \otimes_\Z \Z/b\Z \to \Z/d\Z$，使得交换图
\begin{center}
\begin{tikzcd}[row sep=large, column sep=large]
\Z/a\Z \times \Z/b\Z \arrow[r, "\otimes"] \arrow[dr, "B"'] & \Z/a\Z \otimes_\Z \Z/b\Z \arrow[d, "f"] \\
& \Z/d\Z
\end{tikzcd}
\end{center}
成立，即对简单张量有 $f(x \otimes y) = xy \bmod d$.特别地，$f(x \otimes 1) = x \bmod d$，所以 $f$ 是满射.因此 $\Z/a\Z \otimes_\Z \Z/b\Z$ 的阶至少为 $d$，从而其阶恰好为 $d$，$f$ 是同构.
\end{proof}

\begin{example}\label{ex:3-5-collapse}
阿贝尔群 $\Z/3\Z \otimes_\Z \Z/5\Z = 0$.初见张量积的这种``坍缩''（collapsing）现象常让人难以置信，但它有着极其具体的数学涵义：从 $\Z/3\Z \times \Z/5\Z$ 出发到达任何阿贝尔群 $A$ 的每个 $\Z$-双线性映射 $B$ 都恒为零.这很容易直接验证：
\[
  3B(a, b) = B(3a, b) = B(0, b) = 0 \quad\text{且}\quad 5B(a, b) = B(a, 5b) = B(a, 0) = 0,
\]
因此 $B(a, b)$ 被理想 $3\Z + 5\Z = \Z$ 零化.故而 $B(a, b)$ 被 $1$ 零化，即 $B(a, b) = 0$.

在 $\Z/a\Z \otimes_\Z \Z/b\Z$ 中，所有张量都是简单张量：$x \otimes y = xy(1 \otimes 1)$，而 $1 \otimes 1$ 的倍数之和仍然是 $1 \otimes 1$ 的倍数，所以
\[
  \Z/a\Z \otimes_\Z \Z/b\Z = \Z(1 \otimes 1) = \{x \otimes 1 : x \in \Z\}.
\]
\end{example}

请特别注意定理~\ref{thm:cyclic-tensor} 证明中是如何通过双线性映射 $B \colon \Z/a\Z \times \Z/b\Z \to \Z/d\Z$ 和张量积的普遍映射性质构造出线性映射 $f$ 的.要在 $M \otimes_R N$ 上定义一个把简单张量 $m \otimes n$ 送到指定目标的线性映射，\textbf{永远要退一步，先从 $M \times N$ 出发定义一个把 $(m, n)$ 送到该目标的双线性映射，并严格证明它是双线性的}！随后，普遍映射性质就会免费赋予你所期望的线性映射.正如一位匿名学生曾写道：``做张量积习题如果毫无头绪，那就先从 $M \times N$ 出发构造一个精心挑选的双线性映射，因为除此之外你基本上什么也做不了.''

\begin{theorem}\label{thm:ideal-tensor}
设 $I$ 与 $J$ 为环 $R$ 的理想.则存在唯一的 $R$-模同构
\[
  R/I \otimes_R R/J \cong R/(I + J),
\]
满足 $\bar{x} \otimes \bar{y} \mapsto \overline{xy}$.特别地，取 $I = J = 0$，有 $R \otimes_R R \cong R$（对应 $x \otimes y \mapsto xy$）.

当 $R = \Z$ 且 $I, J$ 非零时，这正是定理~\ref{thm:cyclic-tensor}.
\end{theorem}

\begin{proof}
从函数 $R/I \times R/J \to R/(I + J)$ 开始，定义为 $(x \bmod I, y \bmod J) \mapsto xy \bmod (I + J)$.该函数良定义且双线性，因此由张量积普遍映射性质得到线性映射 $f \colon R/I \otimes_R R/J \to R/(I + J)$，满足 $f(\bar{x} \otimes \bar{y}) = \overline{xy}$.

为了构造逆映射，考虑映射 $R \to R/I \otimes_R R/J$（定义为 $r \mapsto r(\bar{1} \otimes \bar{1})$）.这是 $R$-线性的；当 $r \in I$ 时，其值为 $\bar{r} \otimes \bar{1} = \bar{0} \otimes \bar{1} = 0$.同理当 $r \in J$ 时其值亦为 $0$.因此其核包含 $I + J$，从而诱导出线性映射 $g \colon R/(I + J) \to R/I \otimes_R R/J$，满足 $g(r \bmod (I + J)) = r(\bar{1} \otimes \bar{1}) = \bar{r} \otimes \bar{1} = \bar{1} \otimes \bar{r}$.

直接验证复合：$f(g(r \bmod (I + J))) = f(\bar{r} \otimes \bar{1}) = r \bmod (I + J)$.另一方面，在 $R/I \otimes_R R/J$ 中，简单张量为 $\bar{x} \otimes \bar{y} = xy(\bar{1} \otimes \bar{1})$，故所有张量皆为 $\bar{1} \otimes \bar{1}$ 的标量倍.直接计算即知 $g(f(t)) = t$ 对所有张量 $t$ 成立.因此 $f$ 与 $g$ 互为逆同构.
\end{proof}

\begin{remark}\label{rem:ideal-ops}
对于理想 $I$ 与 $J$，常见的代数运算有 $I + J$、$I \cap J$ 与 $IJ$.交 $I \cap J$ 是线性映射 $R \to R/I \oplus R/J$（定义为 $r \mapsto (\bar{r}, \bar{r})$）的核.而定理~\ref{thm:ideal-tensor} 告诉我们，和 $I + J$ 正好是线性映射 $R \to R/I \otimes_R R/J$（定义为 $r \mapsto r(\bar{1} \otimes \bar{1})$）的核！
\end{remark}

\begin{theorem}\label{thm:quotient-tensor-m}
设 $I$ 为 $R$ 的理想，$M$ 为 $R$-模.则存在唯一的 $R$-模同构
\[
  (R/I) \otimes_R M \cong M/IM,
\]
满足 $\bar{r} \otimes m \mapsto \overline{rm}$.特别地，取 $I = (0)$，有 $R \otimes_R M \cong M$（对应 $r \otimes m \mapsto rm$）；取 $M = R$，有 $R \otimes_R R \cong R$（对应 $r \otimes r' \mapsto rr'$）.
\end{theorem}

\begin{proof}
从双线性映射 $(R/I) \times M \to M/IM$（定义为 $(\bar{r}, m) \mapsto \overline{rm}$）出发，由普遍映射性质诱导出线性映射 $f \colon (R/I) \otimes_R M \to M/IM$，满足 $f(\bar{r} \otimes m) = \overline{rm}$.

逆映射由 $M \to (R/I) \otimes_R M$（定义为 $m \mapsto \bar{1} \otimes m$）诱导：该映射在 $M$ 上是线性的，并且零化 $IM$（因为对 $i \in I$，$\bar{1} \otimes im = \bar{i} \otimes m = \bar{0} \otimes m = 0$），故诱导出线性映射 $g \colon M/IM \to (R/I) \otimes_R M$，满足 $g(\bar{m}) = \bar{1} \otimes m$.容易验证 $f \circ g = \id$ 且 $g \circ f = \id$.
\end{proof}

\begin{example}\label{ex:an-a}
对任意阿贝尔群 $A$，$(\Z/n\Z) \otimes_\Z A \cong A/nA$（通过 $\bar{m} \otimes a \mapsto \overline{ma}$）.
\end{example}

\begin{remark}\label{rem:r-tensor-m}
同构 $R \otimes_R M \cong M$（通过 $r \otimes m \mapsto rm$）表明：从 $R \times M$ 出发的 $R$-双线性映射 $B$ 与从 $M$ 出发的 $R$-线性映射是一一对应的.因为由双线性有 $B(r, m) = B(1, rm)$，而 $B(1, -)$ 在第二个变元上是线性的.
\end{remark}

\begin{remark}\label{rem:ideal-mult-map}
对于理想 $I \subset R$ 与模 $M$，存在典范满同态 $I \otimes_R M \to IM$（满足 $i \otimes m \mapsto im$）.但注意它\textbf{通常不是单射}！例如当 $R = \Z, I = a\Z, M = \Z/a\Z$（$a \geq 2$）时，$I \otimes_R M = (a\Z) \otimes_\Z (\Z/a\Z) \cong \Z/a\Z \neq 0$，然而 $IM = a(\Z/a\Z) = \{0\}$.
\end{remark}

\begin{theorem}\label{thm:free-tensor-basis}
若 $F$ 与 $F'$ 为自由 $R$-模，分别具有基 $\{e_i\}_{i \in I}$ 与 $\{e'_j\}_{j \in J}$，则 $F \otimes_R F'$ 是自由 $R$-模，且以 $\{e_i \otimes e'_j\}_{(i, j) \in I \times J}$ 为一组基.
\end{theorem}

\begin{proof}
若 $F$ 或 $F'$ 为 $0$，结论显然成立.以下设两者均非零（故 $R \neq 0$ 且指标集非空）.由定理~\ref{thm:spanning-tensor}，$\{e_i \otimes e'_j\}$ 生成 $F \otimes_R F'$.

为证明线性无关性，设 $\sum_{i, j} c_{ij} (e_i \otimes e'_j) = 0$（非零项仅有限个）.任取固定的指标对 $(i_0, j_0)$.考虑双线性函数
\[
  F \times F' \to R, \quad (v, w) \mapsto v_{i_0} w_{j_0},
\]
其中 $v = \sum v_i e_i$ 与 $w = \sum w_j e'_j$ 是在给定基下的坐标.由普遍映射性质，存在线性函数 $f_0 \colon F \otimes_R F' \to R$，使得 $f_0(v \otimes w) = v_{i_0} w_{j_0}$.特别地，
\[
  f_0(e_{i_0} \otimes e'_{j_0}) = 1, \quad\text{且当 }(i, j) \neq (i_0, j_0)\text{ 时 } f_0(e_i \otimes e'_j) = 0.
\]
将 $f_0$ 作用于关系式 $\sum c_{ij} (e_i \otimes e'_j) = 0$，即得 $c_{i_0 j_0} = 0$.由于 $(i_0, j_0)$ 是任意的，所有系数皆为 $0$.
\end{proof}

\begin{example}\label{ex:vector-spaces-tensor}
设 $K$ 为域，$V$ 与 $W$ 为有限维非零 $K$-向量空间，分别具有基 $\{e_1, \dots, e_m\}$ 与 $\{f_1, \dots, f_n\}$.定理~\ref{thm:free-tensor-basis} 表明 $V \otimes_K W$ 的每个元素均可唯一表示为 $\sum_{i,j} c_{ij} (e_i \otimes f_j)$，其维度为 $mn = (\dim V)(\dim W)$.
\end{example}

\begin{example}\label{ex:nonelementary}
设 $R \neq 0$，$F$ 与 $F'$ 为秩 $\geq 2$ 的自由 $R$-模，基分别为 $\{e_1, \dots, e_m\}$ 与 $\{e'_1, \dots, e'_n\}$.在 $F \otimes_R F'$ 中，张量
\[
  t = e_1 \otimes e'_1 + e_2 \otimes e'_2
\]
绝不是简单张量.

事实上，任何简单张量形如 $\left(\sum_{i=1}^m a_i e_i\right) \otimes \left(\sum_{j=1}^n b_j e'_j\right) = \sum_{i,j} a_i b_j (e_i \otimes e'_j)$.若它等于 $t$，比较基 $\{e_i \otimes e'_j\}$ 的系数可得：
\[
  a_1 b_1 = 1, \quad a_1 b_2 = 0, \quad a_2 b_1 = 0, \quad a_2 b_2 = 1.
\]
由 $a_1 b_1 = 1$ 与 $a_2 b_2 = 1$ 可知 $a_1$ 与 $b_2$ 均可逆，这与 $a_1 b_2 = 0$ 尖锐矛盾！因此 $e_1 \otimes e'_1 + e_2 \otimes e'_2$ 不是简单张量.类似地，$e_1 \otimes e'_2 + e_2 \otimes e'_1$ 也不是简单张量\footnote{当 $R = K$ 为域时，$K^n \otimes_K K^n$ 中的简单张量由二次多项式方程 $c_{ii} c_{jj} = c_{ij} c_{ji}$ 完全刻画（Segre 嵌入），详见 \cite{grone}.}.
\end{example}

\begin{example}\label{ex:poly-tensor}
对于 $R \neq 0$，$R[X] \otimes_R R[Y]$ 是以 $\{X^i \otimes Y^j\}_{i,j \geq 0}$ 为基的自由 $R$-模，故作为 $R$-模有
\[
  R[X] \otimes_R R[Y] \cong R[X, Y], \quad X^i \otimes Y^j \mapsto X^i Y^j.
\]
更一般地，$R[X_1, \dots, X_k] \cong R[X]^{\otimes k}$.普通多项式乘法与张量积的区别，恰如将单变元多项式相乘写成 $f(T)g(T)$ 与二元多项式 $f(X)g(Y)$ 的区别.
\end{example}

\begin{example}\label{ex:mat-tensor}
回到例~\ref{ex:rank-one}，对于 $v, w \in \R^n$，映射 $B(v, w) = vw^\top \in \mathrm{M}_n(\R)$ 是双线性的，因而诱导线性映射 $L \colon \R^n \otimes_\R \R^n \to \mathrm{M}_n(\R)$，满足 $L(v \otimes w) = vw^\top$.在例~\ref{ex:rank-one} 中我们看到秩 $1$ 矩阵的和不再是秩 $1$ 矩阵（$B$ 的像在加法下不封闭）.而在张量积中，
\[
  B(e_1, e_1) + B(e_2, e_2) = L(e_1 \otimes e_1 + e_2 \otimes e_2),
\]
它完全是张量积上线性映射的一个正常取值！此外，$L(e_i \otimes e_j) = E_{ij}$ 将标准基映为矩阵基，因此 $L$ 是向量空间的同构：$\R^n \otimes_\R \R^n \cong \mathrm{M}_n(\R)$.
\end{example}

\begin{theorem}\label{thm:free-power}
若 $F$ 是以 $\{e_i\}_{i \in I}$ 为基的自由 $R$-模，则对 $k \geq 1$，$k$ 重张量幂 $F^{\otimes k}$ 是以 $\{e_{i_1} \otimes \dots \otimes e_{i_k}\}_{(i_1, \dots, i_k) \in I^k}$ 为基的自由 $R$-模.
\end{theorem}

\begin{theorem}\label{thm:m-tensor-free}
若 $M$ 是任意 $R$-模，而 $F$ 是以 $\{e_i\}_{i \in I}$ 为基的自由 $R$-模，则 $M \otimes_R F$ 中的每个元素都可以\textbf{唯一}表示为
\[
  \sum_{i \in I} m_i \otimes e_i,
\]
其中几乎所有 $m_i$ 均为 $0$.换言之，有典范同构 $M \otimes_R F \cong \bigoplus_{i \in I} M$.
\end{theorem}

\begin{proof}
由定理~\ref{thm:spanning-tensor} 与标量前移规则，每个元素均可写成有限和 $\sum_{i \in I} m_i \otimes e_i$.定义映射 $f \colon \bigoplus_{i \in I} M \to M \otimes_R F$ 为 $f((m_i)) = \sum m_i \otimes e_i$.

逆映射构造如下：函数 $M \times F \to \bigoplus_{i \in I} M$ 定义为 $(m, \sum r_i e_i) \mapsto (r_i m)_{i \in I}$.这是双线性的，故诱导线性映射 $g \colon M \otimes_R F \to \bigoplus_{i \in I} M$.容易验证 $f \circ g = \id$ 与 $g \circ f = \id$.直和坐标的唯一性即蕴含了表示的唯一性.
\end{proof}

\begin{example}\label{ex:s-tensor-poly}
对于含微环扩张 $S \supset R$，$S \otimes_R R[X]$ 中的元素具有唯一表示 $\sum_{n \geq 0} s_n \otimes X^n$，故作为 $R$-模有 $S \otimes_R R[X] \cong S[X]$.
\end{example}

\begin{remark}\label{rem:injectivity-caution}
虽然验证恒等式 $f(g(t)) = t$ 只需在简单张量上检验（因加性），但\textbf{绝不能仅凭简单张量上的性质来证明线性映射 $f \colon M \otimes_R N \to P$ 是单射}！即使对所有简单张量都有 $f(m \otimes n) = 0 \implies m \otimes n = 0$，整个映射 $f$ 仍然完全可能不是单射，因为单射性不是加性性质.因此，证明张量积线性映射为同构最稳妥的途径，往往是显式构造逆映射并验证其双向复合为恒等映射.
\end{remark}

\begin{theorem}\label{thm:fg-power-nonzero}
若 $M$ 是非零有限生成 $R$-模，则对任意 $k \geq 1$，张量幂 $M^{\otimes k} \neq 0$.
\end{theorem}

\begin{proof}
必然有 $R \neq 0$.设 $M = R x_1 + \dots + R x_d$（其中生成元个数 $d \geq 1$ 最小）.令 $N = R x_1 + \dots + R x_{d-1}$（若 $d = 1$ 则 $N = 0$），故 $M = N + R x_d$ 且 $x_d \notin N$.令 $I = \{r \in R : r x_d \in N\}$，则 $I$ 是 $R$ 的真理想（$1 \notin I$）.

任取 $m \in M$ 写成 $n + r x_d$（$n \in N, r \in R$），虽然 $r$ 本身不唯一，但 $r \bmod I$ 是完全良定义的.因此映射 $M^k \to R/I$，定义为
\[
  (n_1 + r_1 x_d, \dots, n_k + r_k x_d) \mapsto r_1 \dots r_k \bmod I,
\]
是良定义的多重线性映射.由普遍映射性质，诱导出 $R$-线性映射 $M^{\otimes k} \to R/I$，它将 $x_d \otimes \dots \otimes x_d$ 映为 $1 \bmod I \neq 0$.因此 $M^{\otimes k} \neq 0$.
\end{proof}

\begin{example}\label{ex:qz-powers}
非有限生成模的张量幂可以完全消失：$(\Q/\Z)^{\otimes 2} = 0$（例~\ref{ex:torsion-collapse}）.这一例子的精妙之处在于：$\Q/\Z = \bigcup_{a \geq 1} (\frac{1}{a}\Z/\Z)$ 是循环子群的并，而每个子群自身的张量平方非零：$(\frac{1}{a}\Z/\Z)^{\otimes 2} \cong \frac{1}{a}\Z/\Z \neq 0$.这深刻揭示了双线性映射的特性：$(\frac{1}{5}) \otimes (\frac{1}{5})$ 在子模自身的张量积中非零，但在大模 $(\Q/\Z)^{\otimes 2}$ 中却坍缩为零！这说明简单张量必须置于具体的张量积背景模中来理解.
\end{example}

以下定理讨论当 $R$ 为整环（domain）时的张量积性质：

\begin{theorem}\label{thm:frac-vector-tensor}
设 $R$ 为整环，$K$ 为其分式域，$V$ 为 $K$-向量空间.则存在 $R$-模同构
\[
  K \otimes_R V \cong V, \quad x \otimes v \mapsto xv.
\]
\end{theorem}

\begin{proof}
乘法 $K \times V \to V$ 是 $R$-双线性的，由普遍映射性质诱导 $R$-线性映射 $f \colon K \otimes_R V \to V$（满足 $f(x \otimes v) = xv$）.显然 $f(1 \otimes v) = v$，故 $f$ 满.

为证单射，我们先证明 $K \otimes_R V$ 中的每个张量都是简单张量且第一项可化为 $1$：设 $x = a/b$（$a, b \in R, b \neq 0$），则
\[
  x \otimes v = \frac{a}{b} \otimes v = \frac{1}{b} \otimes av = \frac{1}{b} \otimes \frac{ab}{b}v = \frac{b}{b} \otimes \frac{a}{b}v = 1 \otimes xv.
\]
这里巧妙利用了 $V$ 上的 $K$-数乘将标量在 $\otimes$ 两侧转移.因此每个张量均可写成 $1 \otimes v$ 的形式.若 $f(1 \otimes v) = 0$，则 $v = 0$，故核平凡，$f$ 为同构.
\end{proof}

\begin{example}\label{ex:q-tensor-q}
当 $V = K$ 时，$K \otimes_R K \cong K$.例如 $\Q \otimes_\Z \Q \cong \Q$，以及 $\Q \otimes_\Z \R \cong \R$.
\end{example}

\begin{theorem}\label{thm:submodule-k-tensor}
设 $R$ 为整环，$K$ 为其分式域，$V$ 为 $K$-向量空间.对 $K$ 内部的任意非零 $R$-子模 $M$，均有 $R$-模同构
\[
  M \otimes_R V \cong V, \quad m \otimes v \mapsto mv.
\]
特别地，对 $R$ 的任意非零理想 $I$，均有 $I \otimes_R K \cong K$.
\end{theorem}

\begin{example}\label{ex:ideal-nonprincipal}
设 $R = \Z[\sqrt{10}], K = \Q(\sqrt{10})$.理想 $I = (2, \sqrt{10})$ 是非主理想，故作为 $R$-模 $I \not\cong R$.然而 $I \otimes_R K \cong R \otimes_R K \cong K$.
\end{example}

\begin{theorem}\label{thm:domain-free-tensor-nonzero}
设 $R$ 为整环，$F$ 与 $F'$ 为自由 $R$-模.若 $x \in F$ 与 $x' \in F'$ 均非零，则在 $F \otimes_R F'$ 中恒有
\[
  x \otimes x' \neq 0.
\]
\end{theorem}

\begin{proof}
设 $\{e_i\}$ 与 $\{e'_j\}$ 分别为 $F$ 与 $F'$ 的基.将非零元展开为 $x = \sum a_i e_i$ 与 $x' = \sum a'_j e'_j$.存在非零系数 $a_{i_0}$ 与 $a'_{j_0}$.由于 $R$ 是整环，$a_{i_0} a'_{j_0} \neq 0$.故 $x \otimes x'$ 在基 $\{e_i \otimes e'_j\}$ 下具有非零坐标，因此 $x \otimes x' \neq 0$.
\end{proof}

\begin{remark}\label{rem:domain-counterexample}
若 $R$ 不是整环，定理~\ref{thm:domain-free-tensor-nonzero} 不成立.例如设 $ab = 0$ 且 $a, b \neq 0$，则在 $R \otimes_R R$ 中 $a \otimes b = ab(1 \otimes 1) = 0$.
\end{remark}

\begin{theorem}\label{thm:domain-torsion-behavior}
设 $R$ 为整环，$K$ 为其分式域，$V$ 为 $K$-向量空间.
\begin{enumerate}
  \item 对任意 $R$-模 $M$，存在 $R$-模同构 $V \otimes_R M \cong V \otimes_R (M/M_{\mathrm{tor}})$，其中 $M_{\mathrm{tor}}$ 为 $M$ 的扭子模.
  \item 若 $M$ 是扭模，则 $V \otimes_R M = 0$；若 $M$ 不是扭模且 $V \neq 0$，则 $V \otimes_R M \neq 0$.
  \item 若 $N \subset M$ 使得商模 $M/N$ 是扭模，则自然映射给出同构 $V \otimes_R N \cong V \otimes_R M$（通过 $v \otimes n \mapsto v \otimes n$）.
\end{enumerate}
\end{theorem}

\begin{corollary}\label{cor:domain-zero-tensor}
在定理~\ref{thm:domain-torsion-behavior} 的设定下，在 $V \otimes_R M$ 中 $v \otimes m = 0$ 当且仅当 $v = 0$ 或 $m \in M_{\mathrm{tor}}$.特别地，由 $m \mapsto 1 \otimes m$ 给出的映射 $M \to K \otimes_R M$ 的核恰好为 $M_{\mathrm{tor}}$.
\end{corollary}

\begin{example}\label{ex:dehn-hilbert}
若 $x \in \R^\times$ 且 $y/\pi$ 为无理数，则在 $\R \otimes_\Z (\R/\pi\Z)$ 中 $x \otimes (y \bmod \pi\Z) \neq 0$.这在现代 Dehn 对 Hilbert 第三问题的解答中扮演了关键角色.
\end{example}

% sections/sec5_properties.tex
\section{张量积的一般性质}
\label{sec:properties}

模的直和同构 $M \oplus N \cong N \oplus M$ 以及 $(M \oplus N) \oplus P \cong M \oplus (N \oplus P)$ 在张量积中有着完美的类似物：$M \otimes_R N \cong N \otimes_R M$ 以及 $(M \otimes_R N) \otimes_R P \cong M \otimes_R (N \otimes_R P)$.此外，$\otimes_R$ 对 $\oplus$ 满足分配律：$M \otimes_R (N \oplus P) \cong (M \otimes_R N) \oplus (M \otimes_R P)$.在实际应用中，这些同构具体是如何运作的，往往比它们单纯的存在性更为重要.

\begin{theorem}[交换性]\label{thm:tensor-comm}
存在唯一的 $R$-模同构 $M \otimes_R N \cong N \otimes_R M$，满足 $m \otimes n \mapsto n \otimes m$.
\end{theorem}

\begin{proof}
我们要构造一个把 $m \otimes n$ 映为 $n \otimes m$ 的线性映射 $M \otimes_R N \to N \otimes_R M$.为此，我们退一步，先从 $M \times N$ 出发构造到达目标模 $N \otimes_R M$ 的映射：定义
\[
  M \times N \to N \otimes_R M, \quad (m, n) \mapsto n \otimes m.
\]
由于 $n \otimes m$ 在 $m$ 与 $n$ 上分别是线性的，因此该映射是双线性的.由张量积的普遍映射性质，存在唯一的线性映射 $f \colon M \otimes_R N \to N \otimes_R M$，使得对所有简单张量均有 $f(m \otimes n) = n \otimes m$.

交换 $M$ 与 $N$ 的角色，同理存在唯一的线性映射 $g \colon N \otimes_R M \to M \otimes_R N$，满足 $g(n \otimes m) = m \otimes n$.

在简单张量上显然有 $f(g(n \otimes m)) = f(m \otimes n) = n \otimes m$ 以及 $g(f(m \otimes n)) = m \otimes n$.由于简单张量生成各自的张量积模，由加性即知 $f \circ g = \id$ 且 $g \circ f = \id$.因此 $f$ 为同构.
\end{proof}

\begin{theorem}[结合律]\label{thm:tensor-assoc}
存在唯一的 $R$-模同构 $(M \otimes_R N) \otimes_R P \cong M \otimes_R (N \otimes_R P)$，满足 $(m \otimes n) \otimes p \mapsto m \otimes (n \otimes p)$.
\end{theorem}

\begin{proof}
由定理~\ref{thm:spanning-tensor}，$(M \otimes_R N) \otimes_R P$ 由元素 $(m \otimes n) \otimes p$ 生成，而 $M \otimes_R (N \otimes_R P)$ 由元素 $m \otimes (n \otimes p)$ 生成.因此其上的线性映射由在这些简单张量上的取值唯一决定.

为构造正向线性映射，考虑函数 $M \times N \times P \to M \otimes_R (N \otimes_R P)$，定义为 $(m, n, p) \mapsto m \otimes (n \otimes p)$.对每个固定的 $p \in P$，映射 $b_p(m, n) = m \otimes (n \otimes p)$ 是双线性的，故诱导线性映射 $f_p \colon M \otimes_R N \to M \otimes_R (N \otimes_R P)$，满足 $f_p(m \otimes n) = m \otimes (n \otimes p)$.

进而考虑函数 $(M \otimes_R N) \times P \to M \otimes_R (N \otimes_R P)$（定义为 $(t, p) \mapsto f_p(t)$）.该函数在 $t$ 上是线性的；同时在简单张量 $t = m \otimes n$ 上容易验证它对 $p$ 也是线性的，从而由加性对所有张量 $t$ 均对 $p$ 线性.故它是一个双线性映射！

由张量积普遍映射性质，存在唯一的线性映射 $f \colon (M \otimes_R N) \otimes_R P \to M \otimes_R (N \otimes_R P)$，满足 $f(t \otimes p) = f_p(t)$，特别地 $f((m \otimes n) \otimes p) = m \otimes (n \otimes p)$.反方向映射 $g$ 的构造完全对称，由生成元的线性性即知它们互为逆同构.
\end{proof}

\begin{theorem}[对有限直和的分配律]\label{thm:tensor-distrib-finite}
存在唯一的 $R$-模同构
\[
  M \otimes_R (N \oplus P) \cong (M \otimes_R N) \oplus (M \otimes_R P),
\]
满足 $m \otimes (n, p) \mapsto (m \otimes n, m \otimes p)$.
\end{theorem}

\begin{proof}
我们直接证明 $(M \otimes_R N) \oplus (M \otimes_R P)$ 具备张量积 $M \otimes_R (N \oplus P)$ 的普遍映射性质：

定义双线性映射 $b \colon M \times (N \oplus P) \to (M \otimes_R N) \oplus (M \otimes_R P)$ 为 $b(m, (n, p)) = (m \otimes n, m \otimes p)$.设 $B \colon M \times (N \oplus P) \to Q$ 为任意给定的双线性映射.若交换图
\begin{equation}
\begin{tikzcd}[row sep=large, column sep=large]
M \times (N \oplus P) \arrow[r, "b"] \arrow[dr, "B"'] & (M \otimes_R N) \oplus (M \otimes_R P) \arrow[d, "L"] \\
& Q
\end{tikzcd} \label{eq:5.1}
\end{equation}
成立，则必有 $L(m \otimes n, 0) = B(m, (n, 0))$ 以及 $L(0, m \otimes p) = B(m, (0, p))$.

由于 $(m, n) \mapsto B(m, (n, 0))$ 与 $(m, p) \mapsto B(m, (0, p))$ 均为双线性映射，分别诱导线性映射 $L_1 \colon M \otimes_R N \to Q$ 与 $L_2 \colon M \otimes_R P \to Q$.定义 $L(t_1, t_2) = L_1(t_1) + L_2(t_2)$，则 $L$ 是线性的且使得图~\eqref{eq:5.1} 严格交换.这证明了 $(M \otimes_R N) \oplus (M \otimes_R P)$ 满足普遍映射性质，因而通过唯一的线性映射与 $M \otimes_R (N \oplus P)$ 典范同构.
\end{proof}

\begin{theorem}[对任意直和的分配律]\label{thm:tensor-distrib-arbitrary}
设 $\{N_i\}_{i \in I}$ 为一族 $R$-模.则存在唯一的 $R$-模同构
\[
  M \otimes_R \bigoplus_{i \in I} N_i \cong \bigoplus_{i \in I} (M \otimes_R N_i),
\]
满足 $m \otimes (n_i)_{i \in I} \mapsto (m \otimes n_i)_{i \in I}$.
\end{theorem}

\begin{proof}
证明与定理~\ref{thm:tensor-distrib-finite} 类似：映射 $b(m, (n_i)_{i \in I}) = (m \otimes n_i)_{i \in I}$ 是双线性的.对任意给定的双线性映射 $B \colon M \times (\bigoplus N_i) \to Q$，限制在各分量上的映射诱导出族 $\{L_i \colon M \otimes_R N_i \to Q\}$，从而唯一确定线性映射 $L((t_i)_{i \in I}) = \sum_{i \in I} L_i(t_i)$.因此直和 $\bigoplus_{i \in I} (M \otimes_R N_i)$ 满足普遍映射性质，结论得证.
\end{proof}

\begin{remark}\label{rem:tensor-prod-counterexample}
定理~\ref{thm:tensor-distrib-arbitrary} 的结论对于模的\textbf{直积}（direct product）\textbf{不再成立}！虽然始终存在自然的 $R$-线性映射
\begin{equation}
  M \otimes_R \prod_{i \in I} N_i \to \prod_{i \in I} (M \otimes_R N_i), \quad m \otimes (n_i) \mapsto (m \otimes n_i), \label{eq:5.2}
\end{equation}
但它通常不是同构.例如取 $R = \Z, M = \Q$ 以及 $N_i = \Z/p^i\Z$（$i \geq 1$）.式~\eqref{eq:5.2} 的右侧为零，因为对每个 $i$ 均有 $\Q \otimes_\Z (\Z/p^i\Z) = 0$（例~\ref{ex:torsion-collapse}）；而左侧为 $\Q \otimes_\Z \prod_{i \geq 1} (\Z/p^i\Z)$，由于无限直积 $\prod_{i \geq 1} (\Z/p^i\Z)$ 不是扭群（元素 $(1, 1, 1, \dots)$ 具有无限阶），由定理~\ref{thm:domain-torsion-behavior} 知左侧\textbf{绝不为零}！
\end{remark}

\begin{theorem}[部分变量张量化的技术引理]\label{thm:partial-tensor-lemma}
设 $M_1, \dots, M_k, N$ 为 $R$-模（$k > 2$），并设函数
\[
  \varphi \colon M_1 \times \dots \times M_{k-2} \times M_{k-1} \times M_k \to N
\]
在固定其他坐标时关于最后两个坐标 $M_{k-1}$ 与 $M_k$ 是双线性的.则存在唯一的函数
\[
  \Phi \colon M_1 \times \dots \times M_{k-2} \times (M_{k-1} \otimes_R M_k) \to N,
\]
使其在固定其他坐标时关于最后坐标 $M_{k-1} \otimes_R M_k$ 是线性的，且满足
\begin{equation}
  \Phi(m_1, \dots, m_{k-2}, m_{k-1} \otimes m_k) = \varphi(m_1, \dots, m_{k-2}, m_{k-1}, m_k). \label{eq:5.3}
\end{equation}
若 $\varphi$ 在所有坐标上是多重线性的，则 $\Phi$ 亦为多重线性的.
\end{theorem}

\begin{proof}
固定 $m_1, \dots, m_{k-2}$ 时，$\varphi$ 限制在 $M_{k-1} \times M_k$ 上是双线性的，由张量积普遍映射性质诱导出唯一的线性映射 $\Phi_{m_1, \dots, m_{k-2}} \colon M_{k-1} \otimes_R M_k \to N$.定义 $\Phi(m_1, \dots, m_{k-2}, t) = \Phi_{m_1, \dots, m_{k-2}}(t)$ 即满足全部要求.
\end{proof}

\begin{theorem}[张量-Hom 伴随同构]\label{thm:tensor-hom-adjoint}
设 $M, N, P$ 为 $R$-模.则存在典范 $R$-模同构
\[
  \operatorname{Bil}_R(M, N; P) \cong \Hom_R(M \otimes_R N, P) \cong \Hom_R(M, \Hom_R(N, P)).
\]
\end{theorem}

\begin{proof}
第一个同构 $\operatorname{Bil}_R(M, N; P) \cong \Hom_R(M \otimes_R N, P)$ 纯粹是张量积普遍映射性质的重述：每个双线性映射 $B$ 一对一对应唯一的线性映射 $L_B$.

对于第二个同构：给定双线性映射 $B \colon M \times N \to P$，对每个 $m \in M$，$B(m, -)$ 是从 $N$ 到 $P$ 的线性映射，故 $m \mapsto B(m, -)$ 构成了线性映射 $f_B \colon M \to \Hom_R(N, P)$.反之，给定线性映射 $L \colon M \to \Hom_R(N, P)$，由 $(m, n) \mapsto (L(m))(n)$ 定义的映射 $B_L \colon M \times N \to P$ 是双线性的.容易检验两方向互为逆同构.
\end{proof}

若记函子 $\mathbf{F}_N(M) = M \otimes_R N$ 以及 $\mathbf{G}_N(P) = \Hom_R(N, P)$，定理~\ref{thm:tensor-hom-adjoint} 写为
\[
  \Hom_R(\mathbf{F}_N(M), P) \cong \Hom_R(M, \mathbf{G}_N(P)).
\]
这完全类似于内积中矩阵与其转置的关系 $Av \cdot w = v \cdot A^\top w$.在范畴论中，$\mathbf{F}_N$ 与 $\mathbf{G}_N$ 称为\newterm{adjoint-functors}{伴随函子对}{adjoint functors}（$\mathbf{F}_N$ 左伴随于 $\mathbf{G}_N$）.

\begin{corollary}\label{cor:dual-tensor}
设 $M, N$ 为 $R$-模.则有 $R$-模同构
\[
  \operatorname{Bil}_R(M, N; R) \cong (M \otimes_R N)^\vee \cong \Hom_R(M, N^\vee) \cong \Hom_R(N, M^\vee).
\]
\end{corollary}

\begin{theorem}\label{thm:dual-tensor-hom}
设 $M$ 与 $N$ 为 $R$-模.存在自然的 $R$-线性映射
\[
  M^\vee \otimes_R N \to \Hom_R(M, N), \quad \varphi \otimes n \mapsto [m \mapsto \varphi(m)n].
\]
若 $M$ 与 $N$ 为有限生成自由模，则该映射为同构.特别地，若 $F$ 为有限自由模，则
\[
  F^\vee \otimes_R F \cong \End_R(F).
\]
\end{theorem}

\begin{proof}
映射 $(\varphi, m, n) \mapsto \varphi(m)n$ 是三线性的，由定理~\ref{thm:partial-tensor-lemma} 诱导出双线性映射 $(M^\vee \otimes_R N) \times M \to N$，进而给出所需的线性映射 $f \colon M^\vee \otimes_R N \to \Hom_R(M, N)$.

当 $M$ 与 $N$ 分别具有有限基 $\{e_j\}$ 与 $\{e'_i\}$ 时，元素 $e_j^\vee \otimes e'_i$ 对应的线性算子将基元素 $e_j$ 映为 $e'_i$，而将其余基向量映为 $0$.在矩阵表示下，这正是初等矩阵 $E_{ij}$.由于所有初等矩阵构成了矩阵空间的一组基，因此 $f$ 将基映为基，必定是同构.
\end{proof}

\begin{remark}\label{rem:dual-tensor-failure}
当模不是有限自由时，定理~\ref{thm:dual-tensor-hom} 中的映射可能既非单射亦非满射.例如取 $R = \Z/p^2\Z$ 以及 $M = N = R/pR$.此时 $M^\vee \cong M$，$M^\vee \otimes_R M \cong M$ 且 $\Hom_R(M, M) \cong M$，但自然映射 $M^\vee \otimes_R M \to \Hom_R(M, M)$ 却是恒等于零的！
\end{remark}

\begin{example}\label{ex:vector-space-dual-tensor}
对有限维 $K$-向量空间 $V$ 与 $W$，有典范同构 $V^\vee \otimes_K W \cong \Hom_K(V, W)$ 以及 $W \otimes_K V^\vee \cong \Hom_K(V, W)$.此时，任何线性映射 $L \in \Hom_K(V, W)$ 都可以无坐标地写为若干简单张量 $\varphi_i \otimes w_i$ 的和.

当 $M$ 与 $N$ 为有限自由模时，利用双对偶同构 $M \cong M^{\vee\vee}$，可得张量积的无坐标描述：
\begin{equation}
  M \otimes_R N \cong \operatorname{Bil}_R(M^\vee, N^\vee; R), \label{eq:5.4}
\end{equation}
其中简单张量 $m \otimes n$ 作用在对偶元对 $(\varphi, \psi)$ 上给出标量 $\varphi(m)\psi(n)$.许多教材（如 Alekseevskij 等人 \cite[第~65~页]{alekseevskij}、O'Neill \cite[第~35~页]{oneill} 以及 Halmos \cite[第~40~页]{halmos}）正是采用式~\eqref{eq:5.4} 作为有限维向量空间张量积的初始定义.
\end{example}

\begin{theorem}[张量秩与线性映射的秩]\label{thm:tensor-rank}
设 $V, W$ 为有限维 $K$-向量空间.将线性映射 $L \colon V \to W$ 视为 $V^\vee \otimes_K W$ 中的张量，则 $L$ 作为线性映射的秩 $\operatorname{rank}(L) = \dim_K L(V)$，恰好等于将其表示为非零简单张量之和时所需的最少项数：
\[
  \operatorname{rank}(L) = \min\left\{d \geq 0 : L = \sum_{i=1}^d \varphi_i \otimes w_i\right\}.
\]
\end{theorem}

\begin{proof}
设 $r = \operatorname{rank}(L)$.若 $L = \sum_{i=1}^d \varphi_i \otimes w_i$，则显然 $L(V) \subset \operatorname{Span}(w_1, \dots, w_d)$，故 $r \leq d$.

反之，选取像空间 $L(V)$ 的一组基 $\{e_1, \dots, e_r\}$.对任意 $v \in V$，将 $L(v)$ 展开为 $\sum_{i=1}^r c_i(v) e_i$.由基展开的唯一性与 $L$ 的线性性，每个坐标系数 $c_i \colon V \to K$ 都是线性的，从而 $c_i \in V^\vee$.因此在 $V^\vee \otimes_K W$ 中有 $L = \sum_{i=1}^r c_i \otimes e_i$，即 $L$ 确可写成 $r$ 个简单张量之和.
\end{proof}

\begin{example}\label{ex:r3-r4-rank}
在 $\R^3 \otimes_\R \R^4$ 中，考虑由标准基展开的 12 项简单张量和：
\begin{align*}
  t &= e_1 \otimes e'_1 + 4e_1 \otimes e'_2 + 2e_1 \otimes e'_3 + 3e_1 \otimes e'_4 + 2e_2 \otimes e'_1 + 5e_2 \otimes e'_2 \\
  &\quad + 3e_2 \otimes e'_3 + 4e_2 \otimes e'_4 + e_3 \otimes e'_1 + 7e_3 \otimes e'_2 + 3e_3 \otimes e'_3 + 5e_3 \otimes e'_4.
\end{align*}
通过同构 $\R^3 \otimes_\R \R^4 \cong \Hom(\R^3, \R^4)$，它对应于 $4 \times 3$ 矩阵
\[
  A = \begin{pmatrix}
    1 & 2 & 1 \\
    4 & 5 & 7 \\
    2 & 3 & 3 \\
    3 & 4 & 5
  \end{pmatrix}.
\]
对 $A$ 进行行化简得到行阶梯形矩阵，发现其秩为 $2$（前两列 $w_1, w_2$ 线性无关，第三列 $w_3 = 3w_1 - w_2$）.由定理~\ref{thm:tensor-rank}，张量 $t$ 的张量秩恰好为 $2$！利用定理~\ref{thm:tensor-rank} 的证明方法，可将 $t$ 精确分解为仅仅两项简单张量之和：
\[
  t = (e_1 + 3e_3) \otimes (e'_1 + 4e'_2 + 2e'_3 + 3e'_4) + (e_2 - e_3) \otimes (2e'_1 + 5e'_2 + 3e'_3 + 4e'_4).
\]
\end{example}

\begin{example}\label{ex:id-tensor-rank}
设 $V$ 为 $n$ 维向量空间，基为 $e_1, \dots, e_n$.则张量 $\sum_{i=1}^n e_i \otimes e_i \in V \otimes_K V$ 的张量秩恰为 $n$，它绝不可能表示为少于 $n$ 个简单张量的和.因为在 $V \otimes_K V \cong \End_K(V)$ 下它对应于恒等算子 $\id_V$，而 $\id_V$ 的秩为 $n$.
\end{example}

\begin{theorem}\label{thm:injective-elementary}
线性映射 $f \colon M \otimes_R N \to P$ 在所有简单张量上是单射，当且仅当 $\ker f$ 中不含任何张量秩为 $1$ 或 $2$ 的非零张量.
\end{theorem}

\begin{proof}
若 $f$ 在简单张量上单射，且 $t \in \ker f$ 的秩至多为 $2$，则可写为 $t = m \otimes n - m' \otimes n'$（若秩 $\leq 1$ 则取 $m' = 0$）.由 $f(t) = 0$ 得 $f(m \otimes n) = f(m' \otimes n')$，单射性给出 $m \otimes n = m' \otimes n'$，故 $t = 0$.

反之，若 $f(m \otimes n) = f(m' \otimes n')$，则差 $m \otimes n - m' \otimes n' \in \ker f$ 的秩至多为 $2$.由假设此差必为 $0$，故 $m \otimes n = m' \otimes n'$.
\end{proof}

\begin{corollary}\label{cor:injective-elementary-not-whole}
线性映射 $f \colon M \otimes_R N \to P$ 在简单张量上单射但在全体张量上不是单射，当且仅当 $\ker f \neq 0$ 且 $\ker f$ 中的所有非零张量的张量秩均至少为 $3$.
\end{corollary}

\begin{example}\label{ex:grinberg-counterexample}
设 $\dim_K V = n \geq 3$.张量 $t = \sum_{i=1}^n e_i \otimes e_i \in V \otimes_K V$ 的秩为 $n \geq 3$.自然商映射
\[
  f \colon V \otimes_K V \to (V \otimes_K V) / Kt
\]
的核为一维子空间 $Kt$.其所有非零元的张量秩均为 $n \geq 3$.由推论~\ref{cor:injective-elementary-not-whole}，$f$ 在所有简单张量上严格单射，但在全体张量上绝非单射！这给出了注~\ref{rem:injectivity-caution} 所警示现象的一个极为精彩的构造.
\end{example}

% sections/sec6_base_ext.tex
\section{基变换}
\label{sec:base-ext}

在代数学中，我们经常需要将一个环上的模替换为另一个相关联环上的模.例如在高等代数中，常常将 $\R^n$ 扩充为 $\C^n$，通过将实坐标扩张为复坐标而构造出一个复向量空间；在多项式理论中，证明 $\Z[X]$ 中多项式的不可约性时，常将其视为 $\Q[X]$ 中的多项式，或者通过将系数模 $p$ 约化而将其置于 $(\Z/p\Z)[X]$ 中来考察.我们将看到，所有这些将模的标量环予以更换的过程（$\R^n \to \C^n$、$\Z[X] \to \Q[X]$、$\Z[X] \to (\Z/p\Z)[X]$），均可以通过张量积以统一的方式精确描述.

设 $f \colon R \to S$ 是交换环的同态.我们可以利用 $f$ 将任意 $S$-模 $N$ 视为 $R$-模，其标量作用定义为 $rn := f(r)n$.特别地，$S$ 自身通过 $rs := f(r)s$ 成为一个 $R$-模.这种将 $S$-模 $N$ 转化为 $R$-模的过程称为\newterm{restriction-scalars}{标量限制}{restriction of scalars}.

\begin{example}\label{ex:subring-rest}
若 $R \subset S$ 且 $f$ 为包含映射（如 $\R \hookrightarrow \C$ 或 $\Q \hookrightarrow \C$），这正是将复向量空间视为实向量空间或有理向量空间的标准方式.
\end{example}

\begin{example}\label{ex:quotient-rest}
若 $S = R/I$ 且 $f$ 为模 $I$ 的自然商同态，则每个 $R/I$-模均可通过令 $r$ 以 $r \bmod I$ 的方式作用而自然成为 $R$-模.
\end{example}

\begin{theorem}\label{thm:s-lin-is-r-lin}
设 $N$ 与 $N'$ 为 $S$-模.若映射 $N \to N'$ 是 $S$-线性的，则当我们把 $N$ 与 $N'$ 视为 $R$-模时，它也是 $R$-线性的.
\end{theorem}

\begin{proof}
设 $\varphi \colon N \to N'$ 为 $S$-线性的，则对所有 $s \in S$ 与 $n \in N$ 均有 $\varphi(sn) = s\varphi(n)$.对任意 $r \in R$：
\[
  \varphi(rn) = \varphi(f(r)n) = f(r)\varphi(n) = r\varphi(n),
\]
因此 $\varphi$ 是 $R$-线性的.
\end{proof}

为了记号清晰，由于我们需要频繁在 $R$-模与 $S$-模之间切换，后文通常用 $M$（或 $M'$ 等）表示原生的 $R$-模，而用 $N$（或 $N'$ 等）表示原生的 $S$-模.由于 $N$ 通过标量限制也是 $R$-模，我们可以形成张量积 $R$-模 $M \otimes_R N$，其中标量可以在 $\otimes$ 两侧转移：
\[
  r(m \otimes n) = (rm) \otimes n = m \otimes rn = m \otimes f(r)n.
\]

\newterm{base-extension-concept}{基变换}{base extension}（或称标量扩张）的目标是逆转标量限制的过程：对于一个给定的 $R$-模 $M$，我们希望构造出一个由形如 $sm$ 的元素组成的 $S$-模，使得当 $s = f(r)$ 时与原有的 $rm$ 严格吻合.这个新的 $S$-模将通过在 $R$-模 $S \otimes_R M$ 上赋予特定的 $S$-模结构来实现.

\begin{theorem}\label{thm:s-tensor-m-struct}
加法群 $S \otimes_R M$ 上存在唯一的 $S$-模结构，使得对所有简单张量满足
\begin{equation}
  s'(s \otimes m) := (s's) \otimes m. \label{eq:6.1}
\end{equation}
该结构与原有的 $R$-模结构相容，即对所有 $r \in R$ 与 $t \in S \otimes_R M$ 均有 $rt = f(r)t$.
\end{theorem}

\begin{proof}
由于简单张量在加法上生成整个张量积，规则~\eqref{eq:6.1} 完全确定了任意 $s' \in S$ 在所有张量上的作用.

为证良定义性，固定 $s' \in S$，考虑映射 $S \times M \to S \otimes_R M$（定义为 $(s, m) \mapsto (s's) \otimes m$）.该映射是 $R$-双线性的，由普遍映射性质诱导出唯一的 $R$-线性映射 $\mu_{s'} \colon S \otimes_R M \to S \otimes_R M$，使得 $\mu_{s'}(s \otimes m) = (s's) \otimes m$.定义 $S$-数乘为 $s't := \mu_{s'}(t)$.容易验证其满足模的各项公理：
\begin{gather*}
  1 \cdot t = t, \quad s(t_1 + t_2) = st_1 + st_2, \\
  (s_1 + s_2)t = s_1 t + s_2 t, \quad s_1(s_2 t) = (s_1 s_2)t.
\end{gather*}
相容性来自：$f(r)(s \otimes m) = (f(r)s) \otimes m = (rs) \otimes m = r(s \otimes m)$.
\end{proof}

完全类似地，右侧张量积 $M \otimes_R S$ 上亦具有唯一的 $S$-模结构，满足 $s'(m \otimes s) = m \otimes (s's)$.由定理~\ref{thm:tensor-comm}，两模通过 $s \otimes m \mapsto m \otimes s$ 给出典范的 $S$-模同构：$S \otimes_R M \cong M \otimes_R S$.

\begin{example}\label{ex:r-mod-i-base-ext}
在定理~\ref{thm:quotient-tensor-m} 中我们证明了 $(R/I) \otimes_R M \cong M/IM$ 作为 $R$-模同构.由于两边均自然具有 $R/I$-模结构，容易验证该同构同时也是 $R/I$-模同构.
\end{example}

\begin{example}\label{ex:s-mod-j-base-ext}
令 $N = S \otimes_R M$.对 $S$ 的任意理想 $J$，$JN$ 是 $N$ 的 $S$-子模，且商模 $N/JN$ 是 $S/J$-模.有典范的 $S/J$-模同构
\[
  (S/J) \otimes_R M \cong N/JN = (S \otimes_R M) / J(S \otimes_R M).
\]
特别地，取 $S = R[X]$ 且 $J = XR[X]$，则 $S/J \cong R$.若 $N = R[X] \otimes_R M$，则商模 $N/XN \cong R \otimes_R M \cong M$.这表明我们可以从标量扩张模 $N$ 唯一直观地恢复出原模 $M \cong N/XN$.
\end{example}

\begin{theorem}\label{thm:free-base-ext}
若 $F$ 是具有基 $\{e_i\}_{i \in I}$ 的自由 $R$-模，则标量扩张 $S \otimes_R F$ 是以 $\{1 \otimes e_i\}_{i \in I}$ 为基的自由 $S$-模.
\end{theorem}

\begin{proof}
由定理~\ref{thm:m-tensor-free}，每个元素均可唯一表示为 $\sum s_i \otimes e_i = \sum s_i (1 \otimes e_i)$，这恰好说明 $\{1 \otimes e_i\}$ 是自由 $S$-基.
\end{proof}

\begin{example}\label{ex:rn-cn-ext}
由定理~\ref{thm:free-base-ext}，作为 $S$-模有 $S \otimes_R R^n \cong S^n$.特别地，
\[
  \C \otimes_\R \R^n \cong \C^n, \quad \C \otimes_\R \mathrm{M}_n(\R) \cong \mathrm{M}_n(\C)
\]
作为复向量空间同构.
\end{example}

\begin{example}\label{ex:poly-base-ext}
作为 $S$-模有 $S \otimes_R R[X] \cong S[X]$.例如：
\[
  \C \otimes_\R \R[X] \cong \C[X], \quad \Q \otimes_\Z \Z[X] \cong \Q[X], \quad (\Z/p\Z) \otimes_\Z \Z[X] \cong (\Z/p\Z)[X].
\]
\end{example}

\begin{example}\label{ex:cn-restrict-extend}
如果我们将 $\C^n$ 视为实向量空间（限制标量，此时 $\dim_\R \C^n = 2n$），然后再将其标量扩张回 $\C$：
\[
  \C \otimes_\R \C^n \cong \C \otimes_\R \R^{2n} \cong \C^{2n} \cong \C^n \oplus \C^n.
\]
这绝非回到原空间 $\C^n$，而是变成了原空间的两个副本！原因在于：一旦限制标量，$\C^n$ 就忘掉了自己原本是复向量空间；当再做标量扩张时，它只是将原本独立的实维度重新复化.
\end{example}

\begin{theorem}\label{thm:direct-sum-base-ext}
对任意一族 $R$-模 $\{M_i\}_{i \in I}$，存在典范 $S$-模同构
\[
  S \otimes_R \bigoplus_{i \in I} M_i \cong \bigoplus_{i \in I} (S \otimes_R M_i).
\]
\end{theorem}

\begin{theorem}\label{thm:domain-rank}
设 $R$ 为整环，$K$ 为其分式域.对任意有限生成 $R$-模 $M$，$K \otimes_R M$ 是有限维 $K$-向量空间，且
\[
  \operatorname{rank}(M) := \dim_K(K \otimes_R M)
\]
等于 $M$ 中 $R$-线性无关元的最大个数，并且是 $M$ 的任意生成集元素个数的下界.
\end{theorem}

\begin{example}\label{ex:ideal-rank-one}
整环 $R$ 的任意非零理想 $I$ 均具有秩 $1$：任取两非零元 $a, b \in I$，有 $ba - ab = 0$，故任意两元均线性相关；同时 $K \otimes_R I \cong K$，故 $\dim_K(K \otimes_R I) = 1$.
\end{example}

\begin{theorem}\label{thm:m-tensor-n-s-mod}
设 $M$ 为 $R$-模，$N$ 为 $S$-模.
\begin{enumerate}
  \item 加法群 $M \otimes_R N$ 上具有唯一的 $S$-模结构，使得 $s(m \otimes_R n) = m \otimes_R sn$.
  \item 存在典范的 $S$-模同构
  \begin{equation}
    M \otimes_R N \cong (S \otimes_R M) \otimes_S N, \quad m \otimes_R n \mapsto (1 \otimes_R m) \otimes_S n. \label{eq:6.3}
  \end{equation}
\end{enumerate}
\end{theorem}

式~\eqref{eq:6.3} 表明：跨越不同环的张量积 $M \otimes_R N$ 可以通过先对 $M$ 作标量扩张 $S \otimes_R M$，随后在 $S$ 上作标准的模张量积来完全实现.

\begin{corollary}\label{cor:iso-base-ext}
若 $M \cong M'$ 为 $R$-模同构，且 $N$ 为 $S$-模，则作为 $S$-模有
\[
  M \otimes_R N \cong M' \otimes_R N \quad\text{以及}\quad N \otimes_R M \cong N \otimes_R M'.
\]
\end{corollary}

\begin{example}\label{ex:invariance-rank-free}
利用张量积可以极其漂亮地证明自由模秩的不变性：设 $R \neq 0$ 且 $R^m \cong R^n$.任取 $R$ 的极大理想 $\mathfrak{m}$，由推论~\ref{cor:iso-base-ext}，在域 $R/\mathfrak{m}$ 上有向量空间同构：
\[
  (R/\mathfrak{m}) \otimes_R R^m \cong (R/\mathfrak{m}) \otimes_R R^n \implies (R/\mathfrak{m})^m \cong (R/\mathfrak{m})^n.
\]
由线性代数中域上向量空间基维数的唯一性，即得 $m = n$！
\end{example}

\begin{theorem}\label{thm:k-vector-tensor-same}
设 $R$ 为整环，$K$ 为其分式域.对任意 $K$-向量空间 $V$ 与 $W$，$V \otimes_R W$ 上通过 $V$ 引入的 $K$-向量空间结构与通过 $W$ 引入的 $K$-向量空间结构\textbf{完全相同}，即对所有 $x \in K$ 恒有
\[
  xv \otimes_R w = v \otimes_R xw.
\]
\end{theorem}

\begin{proof}
设 $x = a/b$（$a, b \in R, b \neq 0$）.利用标量在 $\otimes_R$ 两侧的转移与 $V, W$ 上的 $K$-线性运算：
\[
  \left(\frac{a}{b}v\right) \otimes_R w = \frac{1}{b}v \otimes_R aw = \frac{1}{b}v \otimes_R b\left(\frac{a}{b}w\right) = v \otimes_R \frac{a}{b}w.
\]
\end{proof}

\begin{corollary}\label{cor:v-tensor-r-w-k}
在定理~\ref{thm:k-vector-tensor-same} 的设定下，有典范的 $K$-向量空间同构
\[
  V \otimes_R W \cong V \otimes_K W, \quad v \otimes_R w \mapsto v \otimes_K w.
\]
\end{corollary}

\begin{example}\label{ex:r-tensor-z-r}
由于 $\R$ 是 $\Q$-向量空间，有 $\R \otimes_\Z \R \cong \R \otimes_\Q \R$.注意，虽然通过乘法有满射 $\R \otimes_\Q \R \to \R$（$x \otimes y \mapsto xy$），但它\textbf{绝非单射}（例如 $\pi \otimes 1 - 1 \otimes \pi \neq 0$ 但其像为 $0$）.
\end{example}

\begin{theorem}\label{thm:s-module-assoc-collection}
设 $M, M'$ 为 $R$-模，$N, N'$ 为 $S$-模.
\begin{enumerate}
  \item $M \otimes_R N \cong N \otimes_R M$（作为 $S$-模同构）；
  \item $(M \otimes_R N) \otimes_R M' \cong M \otimes_R (N \otimes_R M')$（作为 $S$-模同构）；
  \item 存在唯一的 $S$-模同构 $(M \otimes_R N) \otimes_S N' \cong M \otimes_R (N \otimes_S N')$；
  \item $N \otimes_R (M \oplus M') \cong (N \otimes_R M) \oplus (N \otimes_R M')$（作为 $S$-模同构）.
\end{enumerate}
\end{theorem}

\begin{corollary}\label{cor:base-ext-of-tensor}
设 $M$ 与 $M'$ 为 $R$-模.存在唯一的 $S$-模同构
\[
  S \otimes_R (M \otimes_R M') \cong (S \otimes_R M) \otimes_S (S \otimes_R M'),
\]
满足 $s \otimes_R (m \otimes_R m') \mapsto s((1 \otimes_R m) \otimes_S (1 \otimes_R m'))$.换言之，\textbf{张量积的标量扩张同构于标量扩张后的模在目标环上的张量积}.
\end{corollary}

\begin{theorem}[基变换的普遍映射性质]\label{thm:base-ext-ump}
设 $M$ 为 $R$-模.对于任意 $S$-模 $N$ 与任意 $R$-线性映射 $\varphi \colon M \to N$，存在唯一的 $S$-线性映射 $\varphi^S \colon S \otimes_R M \to N$ 使得如下交换图成立：
\begin{equation}
\begin{tikzcd}[row sep=large, column sep=large]
M \arrow[r] \arrow[dr, "\varphi"'] & S \otimes_R M \arrow[d, "\exists!\,\varphi^S", dashed] \\
& N
\end{tikzcd}
\end{equation}
其中横向映射为 $m \mapsto 1 \otimes m$，且 $\varphi^S(s \otimes m) = s\varphi(m)$.
\end{theorem}

\begin{theorem}[伴随性：标量扩张左伴随于标量限制]\label{thm:base-ext-adjunction}
设 $M$ 为 $R$-模，$N$ 为 $S$-模.由 $\varphi \mapsto \varphi^S$ 给出的映射诱导了 $S$-模同构：
\begin{equation}
  \Hom_R(M, \operatorname{Res}_{S/R}(N)) \cong \Hom_S(S \otimes_R M, N). \label{eq:6.6}
\end{equation}
\end{theorem}

式~\eqref{eq:6.6} 是现代代数与代数几何中最基础的伴随同构之一，它表明函子 $S \otimes_R -$（标量扩张）是遗忘函子 $\operatorname{Res}_{S/R}$（标量限制）的左伴随函子.

% sections/sec7_physics.tex
\section{物理学中的张量}
\label{sec:physics}

在物理学与工程学中，张量通常不是通过多重线性映射来定义，而是通过它们在不同坐标系下的变换行为来定义的.以下是在许多物理学教科书中常见的一种张量定义：

设 $V$ 为 $n$ 维向量空间（$n \geq 1$）.在 $V$ 上，$0$ 阶张量就是一个标量.对于 $k \geq 1$，$V$ 上的一个 \newterm{contravariant-tensor}{逆变张量}{contravariant tensor}（$k$ 阶）是一个在 $V$ 的每个坐标系下均具有 $n^k$ 个分量的对象 $T$；如果 $\{T^{i_1, \dots, i_k}\}_{1 \leq i_1, \dots, i_k \leq n}$ 与 $\{\widetilde{T}^{i_1, \dots, i_k}\}_{1 \leq i_1, \dots, i_k \leq n}$ 是 $T$ 在 $V$ 的两个坐标系下的分量，则它们满足变换律：
\begin{equation}
  \widetilde{T}^{i_1, \dots, i_k} = \sum_{1 \leq j_1, \dots, j_k \leq n} T^{j_1, \dots, j_k} a_{i_1 j_1} \dots a_{i_k j_k}, \label{eq:7.1}
\end{equation}
其中 $(a_{ij})$ 是将 $V$ 的第一个坐标系用第二个坐标系表示的变换矩阵.简言之，逆变 $k$ 阶张量就是``按规则~\eqref{eq:7.1} 变换的量''.

这里用分量所描述的对象，本质上就是张量幂 $V^{\otimes k}$ 中的一个元素.事实上，所谓``坐标系''就是 $V$ 的一组基.设 $\{e_1, \dots, e_n\}$ 是 $V$ 的一组基，在其中 $T$ 具有分量 $\{T^{i_1, \dots, i_k}\}$，将这些分量作为 $V^{\otimes k}$ 的基 $\{e_{i_1} \otimes \dots \otimes e_{i_k}\}$ 的展开系数：
\[
  \sum_{1 \leq i_1, \dots, i_k \leq n} T^{i_1, \dots, i_k} e_{i_1} \otimes \dots \otimes e_{i_k}.
\]
若将其用另一组基（``新坐标系''）$\{f_1, \dots, f_n\}$ 展开，设 $e_j = \sum_{i=1}^n a_{ij} f_i$，代入并重排求和项：
\begin{align*}
  &\sum_{1 \leq j_1, \dots, j_k \leq n} T^{j_1, \dots, j_k} e_{j_1} \otimes \dots \otimes e_{j_k} \\
  &= \sum_{1 \leq j_1, \dots, j_k \leq n} T^{j_1, \dots, j_k} \left(\sum_{i_1=1}^n a_{i_1 j_1} f_{i_1}\right) \otimes \dots \otimes \left(\sum_{i_k=1}^n a_{i_k j_k} f_{i_k}\right) \\
  &= \sum_{1 \leq i_1, \dots, i_k \leq n} \left(\sum_{1 \leq j_1, \dots, j_k \leq n} T^{j_1, \dots, j_k} a_{i_1 j_1} \dots a_{i_k j_k}\right) f_{i_1} \otimes \dots \otimes f_{i_k} \\
  &= \sum_{1 \leq i_1, \dots, i_k \leq n} \widetilde{T}^{i_1, \dots, i_k} f_{i_1} \otimes \dots \otimes f_{i_k}.
\end{align*}
因此在物理学中，$V$ 上逆变 $k$ 阶张量的分量，正是 $V^{\otimes k}$ 中的元素在 $V^{\otimes k}$ 的某组基下的展开坐标.在物理学中，$V$ 的维数通常为 $3$ 或 $4$.

现在从 $V$ 的张量幂转向其对偶空间 $V^\vee$ 的张量幂.我们要比较 $(V^\vee)^{\otimes \ell}$ 中的元素在由对偶基构成的坐标系下的表示.

为了与物理学和微分几何中的习惯一致，从现在起将基 $\{e_1, \dots, e_n\}$ 与 $\{f_1, \dots, f_n\}$ 的对偶基记作 $\{e^1, \dots, e^n\}$ 与 $\{f^1, \dots, f^n\}$（指标写在上标），满足 $e^i(e_j) = f^i(f_j) = \delta_{ij}$.$V$ 的一般元素写为 $\sum c^i e_i$（基带指标下角标，系数带指标上角标），而对偶空间 $V^\vee$ 的一般元素写为 $\sum c_i e^i$（基带指标上角标，系数带指标下角标）.

任取张量 $\mathsf{T} \in (V^\vee)^{\otimes \ell}$，在基 $\{e^{i_1} \otimes \dots \otimes e^{i_\ell}\}$ 下展开为
\begin{equation}
  \mathsf{T} = \sum_{1 \leq i_1, \dots, i_\ell \leq n} T_{i_1, \dots, i_\ell} e^{i_1} \otimes \dots \otimes e^{i_\ell}. \label{eq:7.2}
\end{equation}
已知在 $V$ 中 $e_j = \sum_{i=1}^n a_{ij} f_i$.计算对偶基的变换关系：
\begin{equation}
  e_j = \sum_{i=1}^n a_{ij} f_i \quad\Longrightarrow\quad f^j = \sum_{i=1}^n a_{ji} e^i. \label{eq:7.3}
\end{equation}
注意式~\eqref{eq:7.3} 中出现了转置矩阵 $(a_{ji})$.将矩阵 $(a_{ij})$ 的逆矩阵记作 $(a^{ij})$，则对偶基的逆向变换为：
\begin{equation}
  e^j = \sum_{i=1}^n a^{ji} f^i. \label{eq:7.6}
\end{equation}
将式~\eqref{eq:7.6} 代入式~\eqref{eq:7.2}，即得分量的协变变换律：
\begin{equation}
  \widetilde{T}_{i_1, \dots, i_\ell} = \sum_{1 \leq j_1, \dots, j_\ell \leq n} T_{j_1, \dots, j_\ell} a^{j_1 i_1} \dots a^{j_\ell i_\ell}. \label{eq:7.7}
\end{equation}
按规则~\eqref{eq:7.7} 变换的量称为 $\ell$ 阶\newterm{covariant-tensor}{协变张量}{covariant tensor}.它无非就是对偶张量空间 $(V^\vee)^{\otimes \ell}$ 中的元素.

更一般地，按如下规则变换的多重指标量：
\begin{equation}
  \widetilde{T}_{j_1 \dots j_\ell}^{i_1 \dots i_k} = \sum_{\substack{1 \leq p_1, \dots, p_k \leq n \\ 1 \leq q_1, \dots, q_\ell \leq n}} T_{q_1 \dots q_\ell}^{p_1 \dots p_k} a_{i_1 p_1} \dots a_{i_k p_k} a^{q_1 j_1} \dots a^{q_\ell j_\ell} \label{eq:7.8}
\end{equation}
称为\newterm{type-k-l-tensor}{$(k, \ell)$ 型张量}{tensor of type (k, l)}（其总阶数为 $k + \ell$）.它正是混合张量空间
\[
  V^{\otimes k} \otimes (V^\vee)^{\otimes \ell}
\]
中的元素.

\begin{example}\label{ex:rank2-tensor-physics}
$2$ 阶张量分为三类：
\begin{enumerate}
  \item \textbf{$(2, 0)$ 型张量}：属于 $V^{\otimes 2}$，其矩阵变换规则为 $(\widetilde{T}^{ij}) = A (T^{ij}) A^\top$；
  \item \textbf{$(0, 2)$ 型张量}：属于 $(V^\vee)^{\otimes 2}$，其矩阵变换规则为 $(\widetilde{T}_{ij}) = (A^{-1})^\top (T_{ij}) A^{-1}$，它对应于 $V$ 上的\textbf{双线性型}（例如内积度规、闵氏度规）；
  \item \textbf{$(1, 1)$ 型张量}：属于 $V \otimes V^\vee \cong \End(V)$，其矩阵变换规则为 $(\widetilde{T}^i_j) = A (T^i_j) A^{-1}$，它对应于 $V$ 上的\textbf{线性变换}（相似变换）.
\end{enumerate}
\end{example}

\begin{remark}\label{rem:einstein-quote}
数学家与物理学家看待张量的角度截然不同：
\begin{itemize}
  \item \textbf{数学家}：张量是张量空间（由多重线性映射普遍性质定义的模或向量空间）中的元素，完全独立于任何基或坐标系.
  \item \textbf{物理学家}：``张量是由一个或多个指标组织的分量体系，在坐标变换下服从特定的变换规则''\cite{feynman-lectures}.
\end{itemize}

Einstein 在其 1916 年关于广义相对论的基础论文中写道 \cite[第~157~页]{einstein-gr}：
\begin{quote}
``设对于任意坐标系，某些对象（'张量'）由若干坐标函数定义，这些函数称为张量的'分量'.当从原坐标系变换到新坐标系且变换已知时，存在确定的法则来计算新坐标系下的这些分量.被称为张量的对象其特征在于其分量的变换方程是线性且齐次的.因此，若所有分量在原坐标系中全为零，则它们在新坐标系中亦全部为零.''
\end{quote}
Einstein 这里所说的``线性且齐次变换''，在现代数学语言中就是多重线性变换.

Hassler Whitney 在回忆早年学习张量的困惑时曾生动地写道 \cite[第~114~页]{whitney-history}：
\begin{quote}
``我不得不摆弄张量；但我连见都不被允许见到它们，只能被告知它们在坐标变换下不断变换的外衣装扮，我又怎能掌握它们呢？我无论如何都必须一把揪出这些淘气的小家伙，直面它们的真面目！''
\end{quote}
\end{remark}

\begin{example}\label{ex:stress-tensor}
在连续介质力学中，最基础的 $2$ 阶张量是柯西应力张量.在一点处施加力时，产生的应力方向未必沿着力的方向，且这种响应在输入力上是线性的.因此某点处的应力由一个线性变换描述，从而是一个 $(1, 1)$ 型张量（属于 $V \otimes V^\vee$）.由于受力情况在物体各点处不同，应力实际上是由各点处张量构成的光滑族，即\newterm{tensor-field}{张量场}{tensor field}.

在微分流形 $M$ 上，张量场可以在局部定义（在每点 $p \in M$ 选取切空间与余切空间张量积 $T_p(M)^{\otimes k} \otimes T_p(M)^{\vee \otimes \ell}$ 中的元素并随 $p$ 光滑变化），亦可在全局定义为光滑切向量场模 $\mathfrak{X}(M)$ 在光滑函数环 $C^\infty(M)$ 上的张量积：
\[
  \mathfrak{X}(M)^{\otimes k} \otimes_{C^\infty(M)} \mathfrak{X}(M)^{\vee \otimes \ell}.
\]
此时局部坐标变换 $x \mapsto y$ 诱导的变换矩阵正是雅可比偏导数矩阵 $\frac{\partial y^i}{\partial x^j}$.
\end{example}

需要注意，\textbf{并非所有带指标的量都是张量}.例如微分几何中的克里斯托费尔符号（Christoffel symbols）$\Gamma^k_{ij}$，虽然带有三个指标，但在坐标变换下包含偏导数的二次非齐次项，不满足多重线性变换法则，因此它不是张量.

在量子力学中，张量积扮演着更为根本的角色.在经典力学中，复合力学系统的相空间是子系统流形的笛卡尔积 $M \times N$；而在量子力学中，系统的状态由复 Hilbert 空间（如 $L^2(\R^3)$）中的单位射线描述.当组合两个量子系统时，复合系统的状态空间是各自 Hilbert 空间的张量积：
\[
  \mathcal{H}_1 \otimes_\C \mathcal{H}_2.
\]
在直积 $M \times N$ 中，每个点都是形如 $(x, y)$ 的平凡对；但在张量积 $\mathcal{H}_1 \otimes_\C \mathcal{H}_2$ 中，绝大多数张量都\textbf{不是简单张量}.那些无法表示为单一简单张量 $v \otimes w$ 的态，正是量子力学中最神奇也最深刻的现象——\newterm{entanglement}{量子纠缠态}{entangled states}.

% sections/bibliography.tex
\begin{thebibliography}{99}
\addcontentsline{toc}{section}{参考文献}

\bibitem{alekseevskij}
D.~V. Alekseevskij, V.~V. Lychagin, A.~M. Vinogradov, \textit{Geometry I}, Springer-Verlag, Berlin, 1991.

\bibitem{bourbaki-alg}
N.~Bourbaki, \textit{Livre II Alg\`ebre Chapitre III (\'etat 4) Alg\`ebre Multilin\'eaire}, 1948. 
\url{http://archives-bourbaki.ahp-numerique.fr/files/original/97a9fed708bdde4dc55547ab5a8ff943.pdf}.

\bibitem{carroll}
S.~Carroll, \textit{Spacetime and Geometry: An Introduction to General Relativity}, Benjamin Cummings, 2003.

\bibitem{einstein-gr}
A.~Einstein, \textit{The Foundation of the General Theory of Relativity} (English translation), pp.~146--199 in \textit{The Collected Papers of Albert Einstein, Vol. 6: The Berlin Years: Writings, 1914--1917}, Princeton Univ. Press, Princeton, 1997. 
\url{https://einsteinpapers.press.princeton.edu/vol6-trans/158}.

\bibitem{eisenbud-harris}
D.~Eisenbud, J.~Harris, \textit{The Geometry of Schemes}, Graduate Texts in Mathematics, vol.~197, Springer-Verlag, New York, 2000.

\bibitem{feynman-lectures}
R.~P. Feynman, Chap.~31 of \textit{The Feynman Lectures on Physics, Vol. II}, Millennium edition, 
\url{http://www.feynmanlectures.caltech.edu/II_31.html}.

\bibitem{gibbs-vector}
J.~W. Gibbs, \textit{Elements of Vector Analysis Arranged for the Use of Students in Physics}, Tuttle, Morehouse \& Taylor, New Haven, 1884. 
\url{https://archive.org/details/elementsvectora00gibb}.

\bibitem{gibbs-grassmann}
J.~W. Gibbs, \textit{On Multiple Algebra}, Proceedings of the American Association for the Advancement of Science, \textbf{35} (1886). 
\url{http://archive.org/details/onmultiplealgeb00gibbgoog}.

\bibitem{grone}
R.~Grone, \textit{Decomposable Tensors as a Quadratic Variety}, Proc. Amer. Math. Soc. \textbf{64} (1977), 227--230.

\bibitem{halmos}
P.~Halmos, \textit{Finite-Dimensional Vector Spaces}, Springer-Verlag, New York, 1974.

\bibitem{hastad}
J.~H\aa stad, \textit{Tensor rank is NP-complete}, J. Algorithms \textbf{11} (1990), 644--654.

\bibitem{hillar-lim}
C.~J. Hillar, L.-H. Lim, \textit{Most tensor problems are NP-hard}, J. ACM \textbf{60} (2013), Art.~45, 39~pp.

\bibitem{kun}
J.~Kun, \textit{How to Conquer Tensorphobia}, 2014. 
\url{https://jeremykun.com/2014/01/17/how-to-conquer-tensorphobia/}.

\bibitem{lee-manifolds}
J.~M. Lee, \textit{Manifolds and Differential Geometry}, Graduate Studies in Mathematics, vol.~107, Amer. Math. Soc., Providence, 2009.

\bibitem{ricci-english}
R.~Hermann, \textit{Ricci and Levi-Civita's Tensor Analysis Paper}, Math Sci Press, Brookline, 1975.

\bibitem{misner-thorne-wheeler}
C.~W. Misner, K.~S. Thorne, J.~A. Wheeler, \textit{Gravitation}, W. H. Freeman and Co., San Francisco, 1973.

\bibitem{murray-vonneumann}
F.~J. Murray, J.~von Neumann, \textit{On Rings of Operators}, Annals of Math. \textbf{37} (1936), 116--229.

\bibitem{oneill}
B.~O'Neill, \textit{Semi-Riemannian Geometry}, Academic Press, New York, 1983.

\bibitem{osgood}
B.~Osgood, Chapter~8 of \textit{The Fourier Transform and its Applications}, 
\url{http://see.stanford.edu/materials/lsoftaee261/chap8.pdf}.

\bibitem{grassmann}
H.-J. Petsche, \textit{Hermann Grassmann -- Biography}, Birkh\"auser, Basel, 2009.

\bibitem{rainich}
G.~Y. Rainich, \textit{Tensor analysis without coordinates}, Proc. Natl. Acad. Sci. USA \textbf{9} (1923), 179--183.

\bibitem{ricci-levicivita}
G.~Ricci, T.~Levi-Civita, \textit{M\'ethodes de Calcul Diff\'erentiel Absolu et Leurs Applications}, Math. Annalen \textbf{54} (1901), 125--201.

\bibitem{shitov}
Y.~Shitov, \textit{How hard is the tensor rank?}, 
\url{https://arxiv.org/abs/1611.01559}.

\bibitem{temple}
G.~Temple, \textit{Cartesian Tensors}, Wiley, New York, 1960.

\bibitem{voigt}
W.~Voigt, \textit{Die fundamentalen physikalischen Eigenschaften der Krystalle in elementarer Darstellung}, Verlag von Veit \& Comp., Leipzig, 1898.

\bibitem{wald}
R.~Wald, \textit{General Relativity}, Univ. Chicago Press, Chicago, 1984.

\bibitem{weyl-groups}
H.~Weyl, \textit{The Theory of Groups and Quantum Mechanics}, E. P. Dutton, New York, 1932.

\bibitem{whitney}
H.~Whitney, \textit{Tensor Products of Abelian Groups}, Duke Math. Journal \textbf{4} (1938), 495--528.

\bibitem{whitney-history}
H.~Whitney, \textit{Moscow 1935: topology moving toward America}, pp.~97--117 in \textit{A Century of Mathematics in America, Part I}, Amer. Math. Soc., Providence, 1988.

\end{thebibliography}


\printterminology

\end{document}
